\exercisesection{\S~4}

The Lie algebras considered in this paragraph are not necessarily finite dimensional.

\begin{enumerate}
\def\labelenumi{\arabic{enumi})}
\tightlist
\item
  Retain the notations of no. 2. Let \(\lambda \in k^{\mathrm{B}}\) . Associate to any \(\alpha \in \mathrm{B}\) the endomorphisms \(X_{-\alpha}^{\lambda}, H_{\alpha}^{\lambda}, X_{\alpha}^{\lambda}\) of the vector space E such that
\end{enumerate}

\[
X _ {- \alpha} ^ {\lambda} (\alpha_ {1}, \ldots , \alpha_ {n}) = (\alpha , \alpha_ {1}, \ldots , \alpha_ {n})
\]

\[
H _ {\alpha} ^ {\lambda} (\alpha_ {1}, \ldots , \alpha_ {n}) = (\lambda (\alpha) - \sum_ {i = 1} ^ {n} n (\alpha_ {i}, \alpha)) (\alpha_ {1}, \ldots , \alpha_ {n}).
\]

The vector \(X_{\alpha}^{\lambda}(\alpha_{1},\ldots,\alpha_{n})\) is defined by induction on n by the formula

\[
X _ {\alpha} ^ {\lambda} (\alpha_ {1}, \ldots , \alpha_ {n}) = (X _ {- \alpha_ {1}} ^ {\alpha} X _ {\alpha} ^ {\lambda} - \delta_ {\alpha , \alpha_ {1}} H _ {\alpha} ^ {\lambda}) (\alpha_ {2}, \ldots , \alpha_ {n}),
\]

where \(\delta_{\alpha, \alpha_1}\) is the Kronecker symbol; we agree that, if \((\alpha_1, \ldots, \alpha_n)\) is the empty word, then \(X_{\alpha}^{\lambda}(\alpha_1, \ldots, \alpha_n)\) is zero.

Show that Lemmas 1 and 2 remain true for these endomorphisms. We thus obtain a representation \(\rho_{\lambda}:\mathfrak{a}\to \mathfrak{gl}(\mathrm{E})\) such that

\[
\rho_ {\lambda} (x _ {\alpha}) = X _ {\alpha} ^ {\lambda}, \rho_ {\lambda} (h _ {\alpha}) = H _ {\alpha} ^ {\lambda}, \rho_ {\lambda} (x _ {- \alpha}) = X _ {- \alpha} ^ {\lambda}.
\]

¶ 2) Retain the notations of no. 3. Let m be an ideal of a.

\begin{enumerate}
\def\labelenumi{\alph{enumi})}
\item
  Let \(\alpha\) and \(\beta\) be two distinct elements of B. Assume that (ad \(x_{\alpha})^{\mathrm{N}}x_{\beta}\) belongs to \(\mathfrak{m}\) for N sufficiently large. Show that \(x_{\alpha \beta} \in \mathfrak{m}\) . (Apply the results of §1, no. 2 to \(\mathfrak{a} / \mathfrak{m}\) , provided with a suitable \(\mathfrak{sl}(2,k)\) -module structure.)
\item
  Show that \(n + \theta n\) is the smallest finite codimensional ideal of a. Show that this is also the smallest ideal containing the \((\operatorname{ad} x_{\alpha})^{4}x_{\beta}\) and the \((\operatorname{ad} x_{-\alpha})^{4}x_{-\beta}\) .
\end{enumerate}

\begin{enumerate}
\def\labelenumi{\arabic{enumi})}
\setcounter{enumi}{2}
\tightlist
\item
  Let \((\mathfrak{g},\mathfrak{h})\) be a split semi-simple Lie algebra, R the corresponding root system, and B a basis of R. For any \(\alpha \in \mathrm{B}\) (resp. for any pair \((\alpha ,\beta)\in \mathrm{B}^2\) ), let \(\mathrm{R}(\alpha)\) (resp. \(\mathrm{R}(\alpha ,\beta)\) ) be the closed subset of R formed by the \(\pm \alpha\) (resp. the smallest closed subset of R containing \(\pm \alpha\) and \(\pm \beta\) ). Let \(\mathfrak{g}(\alpha)\) (resp. \(\mathfrak{g}(\alpha ,\beta)\) ) be the derived algebra of the algebra \(\mathfrak{h} + \mathfrak{g}^{\mathrm{R}(\alpha)}\) (resp. \(\mathfrak{h} + \mathfrak{g}^{\mathrm{R}(\alpha ,\beta)}\) ), cf.~§3.
\end{enumerate}

\begin{enumerate}
\def\labelenumi{\alph{enumi})}
\item
  Show that \(\mathfrak{g}(\alpha) = kH_{\alpha}\oplus \mathfrak{g}^{\alpha}\oplus \mathfrak{g}^{-\alpha}\) ; it is isomorphic to \(\mathfrak{sl}(2,k)\) .
\item
  Show that \(\mathfrak{g}(\alpha,\beta)\) is semi-simple, and that it is generated by \(\mathfrak{g}(\alpha)\) and \(\mathfrak{g}(\beta)\) . Its root system can be identified with \(\mathrm{R}(\alpha,\beta)\) .
\item
  Let \(\mathfrak{s}\) be a Lie algebra (not necessarily finite dimensional). For all \(\alpha \in \mathrm{B}\) , let \(f_{\alpha}\) be a homomorphism from \(\mathfrak{g}(\alpha)\) to \(\mathfrak{s}\) . Assume that, for any pair \((\alpha, \beta)\) , there exists a homomorphism \(f_{\alpha \beta}: \mathfrak{g}(\alpha, \beta) \to \mathfrak{s}\) that extends both \(f_{\alpha}\) and \(f_{\beta}\) . Show that, in that case, there exists a unique homomorphism \(f: \mathfrak{g} \to \mathfrak{s}\) that extends the \(f_{\alpha}\) . (Use Prop. 4 (i).)
\end{enumerate}

\begin{enumerate}
\def\labelenumi{\arabic{enumi})}
\setcounter{enumi}{3}
\item
  Let \(\mathfrak{g}\) be a splittable semi-simple Lie algebra, and \(\sigma\) an automorphism of \(k\) . Let \(\mathfrak{g}_{\sigma}\) be the Lie algebra obtained from \(\mathfrak{g}\) by extending scalars by means of \(\sigma\) . Show that \(\mathfrak{g}_{\sigma}\) is isomorphic to \(\mathfrak{g}\) . (Use the Cor. of Prop. 4.)
\item
  \begin{enumerate}
  \def\labelenumii{\alph{enumii})}
  \tightlist
  \item
    Let g be a simple Lie algebra, and \(k_{1}\) the commutant of the adjoint representation of g. Show that \(k_{1}\) is a commutative field, a finite extension of k, and that g is an absolutely simple \(k_{1}\) -Lie algebra.
  \end{enumerate}
\end{enumerate}

Conversely, if \(k_{1}\) is a finite extension of k, and g an absolutely simple \(k_{1}\) -Lie algebra, then g is a simple k-Lie algebra, and the commutant of its adjoint representation can be identified with \(k_{1}\) .

\begin{enumerate}
\def\labelenumi{\alph{enumi})}
\setcounter{enumi}{1}
\tightlist
\item
  Let \(k'\) be a Galois extension of k containing \(k_{1}\) . Show that \(\mathfrak{g}_{(k')}\) is a product of \([k_{1}:k]\) absolutely simple algebras. When \(\mathfrak{g}_{(k')}\) is split, these algebras are mutually isomorphic. (Use Exerc. 4.)
\end{enumerate}

\begin{enumerate}
\def\labelenumi{\arabic{enumi})}
\setcounter{enumi}{5}
\tightlist
\item
  Let \(\mathbf{A}\) be a commutative ring, and let \(\mathfrak{u}\) be the A-Lie algebra defined by the family of generators \(\{x, y\}\) and the relations
\end{enumerate}

\[
[ x, [ x, y ] ] = 0, \quad [ y, [ y, [ y, x ] ] ] = 0.
\]

Show that \(\mathfrak{u}\) is a free A-module with basis

\[
\{x, y, [ x, y ], [ y, [ x, y ] ] \}.
\]

Show that, when \(\mathrm{A} = k\) , \(\mathfrak{u}\) is isomorphic to the algebra \(\mathfrak{a}_{+} / \mathfrak{n}\) corresponding to a root system of type \(\mathrm{B}_2\) .

¶7) Let A be a commutative ring in which 2 is invertible, and let u be the A-Lie algebra defined by the family of generators \(\{x,y\}\) and the relations

\[
[ x, [ x, y ] ] = 0, \quad [ y, [ y, [ y, [ y, x ] ] ] ] = 0.
\]

Show that \(\mathfrak{u}\) is a free A-module with basis

\[
\left. \right.\left\{x, y, [ x, y ], [ y, [ x, y ] ], [ y, [ y, [ x, y ] ] ], [ x, [ y, [ y, [ x, y ] ] ] ] \right\}.
\]

Show that, when A = k, u is isomorphic to the algebra \(a_{+}/n\) corresponding to a root system of type \(G_{2}\) .

\exercisesection{\S~5}

\begin{enumerate}
\def\labelenumi{\arabic{enumi})}
\item
  The index of \(\operatorname{Aut}_{0}(\mathfrak{g})\) in \(\operatorname{Aut}(\mathfrak{g})\) is finite.
\item
  We have \(\mathrm{Aut}_e(\mathfrak{g}) = \mathrm{Aut}(\mathfrak{g})\) if \(\mathfrak{g}\) is splittable, simple, and of type \(\mathrm{G}_2, \mathrm{F}_4\) or \(\mathrm{E}_8\) .
\item
  Let \(\mathfrak{h}\) be a splitting Cartan subalgebra of \(\mathfrak{g}\) , \(\mathfrak{b}\) a Borel subalgebra of \((\mathfrak{g},\mathfrak{h})\) , \(\mathfrak{n} = [\mathfrak{b},\mathfrak{b}]\) and \(N = \exp \operatorname{ad}_{\mathfrak{g}}\mathfrak{n}\) . Then
\end{enumerate}

\[
\operatorname{Aut} (\mathfrak {g}) = \mathrm{N}. \operatorname{Aut} (\mathfrak {g}, \mathfrak {h}). \mathrm{N}.
\]

(Let \(s \in \operatorname{Aut}(\mathfrak{g})\) . Apply Prop. 10 of §3, no. 3 to \(\mathfrak{b} \cap s(\mathfrak{b})\) , then apply Chap. VII, §3, no. 4, Th. 3.)

\begin{enumerate}
\def\labelenumi{\arabic{enumi})}
\setcounter{enumi}{3}
\tightlist
\item
  Let \(\mathfrak{h}\) be a Cartan subalgebra of \(\mathfrak{g}\) , and \(s\) an element of \(\operatorname{Aut}(\mathfrak{g},\mathfrak{h})\) such that \(sH \neq H\) for all non-zero \(H\) in \(\mathfrak{h}\) . Show that \(s\) is of finite order. (Reduce to the case in which \(\mathfrak{h}\) is splitting, and choose an integer \(n \geq 1\) such that \(\varepsilon(s)^n = 1\) . Then there exists \(\varphi \in \mathrm{T}_{\mathrm{Q}}\) such that \(f(\varphi) = s^n\) . Let \(\sigma\) be the transpose of \(s|\mathfrak{h}\) . Show that \(1 + \sigma + \sigma^2 + \cdots + \sigma^{n-1} = 0\) , and deduce that \(s^{n^2} = 1\) .)
\end{enumerate}

¶ 5) a) Let \(a \in \operatorname{Aut}(\mathfrak{g})\) , and \(\mathfrak{n}\) the nilspace of \(a - 1\) . Show that the following conditions are equivalent:

\begin{enumerate}
\def\labelenumi{(\roman{enumi})}
\item
  \(\operatorname{Ker}(a - 1)\) is a Cartan subalgebra of \(\mathfrak{g}\) .
\item
  \(\mathfrak{n}\) is a Cartan subalgebra of \(\mathfrak{g}\) , and \(a\in \mathrm{Aut}_0(\mathfrak{g})\) .
\item
  \(\dim\mathfrak{n}=\operatorname{rk}(\mathfrak{g})\) and \(a\in\operatorname{Aut}_{0}(\mathfrak{g})\) .
\end{enumerate}

\begin{enumerate}
\def\labelenumi{\alph{enumi})}
\setcounter{enumi}{1}
\item
  Assume from now on that k is algebraically closed. Let V be a vector space, R a root system in V, \(T_{Q}\) the group \(\operatorname{Hom}(Q(R), k^{*})\) , n an integer \(\geq 1\) and \(T_{n}\) the subgroup of \(T_{Q}\) consisting of the elements whose order divides n.~Let \(\zeta\) be a primitive \(n\) th root of unity in \(k\) . For all \(H \in \mathrm{P}(\mathrm{R}^{\vee})\) , let \(\psi(H)\) be the element \(\gamma \mapsto \zeta^{\gamma(H)}\) of \(\mathrm{T}_{\mathrm{Q}}\) . The map \(\psi\) is a homomorphism from \(\mathrm{P}(\mathrm{R}^{\vee})\) to \(\mathrm{T}_n\) , with kernel \(n\mathrm{P}(\mathrm{R}^{\vee})\) . Let \(t \in \mathrm{T}_n\) , and let \(\mathrm{C}\) be an alcove in \(\mathrm{P}(\mathrm{R}^{\vee}) \otimes \mathbf{R}\) . There exists \(w \in \mathrm{W}(\mathrm{R})\) and \(H \in \mathrm{P}(\mathrm{R}^{\vee})\) such that \(\frac{1}{n} H \in \overline{\mathrm{C}}\) and \(\psi(wH) = t\) . (Use Chap. VI, §2, no. 1.)
\item
  Let \(\mathfrak{h}\) be a Cartan subalgebra of \(\mathfrak{g}\) , and \(\mathbb{R}\) the root system of \((\mathfrak{g},\mathfrak{h})\) . Let \(n,\zeta ,\psi\) be as in \(b\) ), and \(f\) as in no. 2. Let \(H\in \mathrm{P}(\mathrm{R}^{\vee})\) . The set of elements of \(\mathfrak{g}\) invariant under \(f(\psi (H))\) is \(\mathfrak{h}\oplus \sum_{\alpha \in \mathbb{R}'}\mathfrak{g}^\alpha\) , where \(\mathbb{R}'\) is the set of \(\alpha \in \mathbb{R}\) such that \(\alpha (H)\in n\mathbf{Z}\) , and \(f(\psi (H))\) satisfies the conditions of \(a\) ) if and only if \(\frac{1}{n} H\) belongs to an alcove.
\end{enumerate}

\(d\) ) Assume from now on that \(\mathfrak{g}\) is simple. Let \(\mathfrak{h}\) and \(\mathbb{R}\) be as in \(c\) ), \((\alpha_{1},\ldots ,\alpha_{l})\) a basis of \(\mathbb{R}\) , \(h\) the Coxeter number of \(\mathbb{R}\) , \(\zeta\) a primitive \(h\) th root of unity in \(k\) , and \(H\) the element of \(\mathfrak{h}\) such that \(\alpha_{i}(H) = 1\) for \(i = 1,\dots,l\) . Prove the following properties:

\begin{enumerate}
\def\labelenumi{(\roman{enumi})}
\item
  The homomorphism \(\gamma \mapsto \zeta^{\gamma (H)}\) from Q(R) to \(k^{*}\) defines an element of \(\operatorname {Aut}(\mathfrak{g},\mathfrak{h})\) which satisfies the conditions in \(a\) ), and has order \(h\) . (Use \(c\) ), Chap. VI, §2, Prop. 5 and Chap. VI, §1, Prop. 31.)
\item
  Every automorphism of \(\mathfrak{g}\) of finite order satisfies the conditions in \(a\) ) and is of order \(\geq h\) .
\item
  The automorphisms of \(\mathfrak{g}\) of order \(h\) satisfying the conditions of \(a\) ) form a conjugacy class in \(\mathrm{Aut}_e(\mathfrak{g})\) . (Use Prop. 5 of no. 3.)
\end{enumerate}

\begin{enumerate}
\def\labelenumi{\alph{enumi})}
\setcounter{enumi}{4}
\tightlist
\item
  Let \(\mathfrak{h}\) and \(\mathrm{R}\) be as in \(c\) ), and let \(w\) be a Coxeter transformation in \(\mathrm{W(R)}\) . Let \(s \in \operatorname{Aut}(\mathfrak{g}, \mathfrak{h})\) be such that \(\varepsilon(s) = w\) . Show that \(s\) satisfies the conditions in \(a\) ), and is of order \(h\) . (Use Chap. VI, §1, Prop. 33, Chap. V, §6, no. 2 and Chap. VII, §4, Prop. 9.)
\end{enumerate}

\(f)\) If \(s \in \operatorname{Aut}(\mathfrak{g})\) , the following conditions are equivalent:

\begin{enumerate}
\def\labelenumi{(\roman{enumi})}
\item
  \(s\) satisfies the conditions in \(a\) ), and is of order \(h\) ;
\item
  there exists a Cartan subalgebra \(\mathfrak{h}\) of \(\mathfrak{g}\) stable under \(s\) such that \(s|\mathfrak{h}\) is a Coxeter transformation in the Weyl group of \((\mathfrak{g},\mathfrak{h})\) . (Use \(d\) ) and \(e\) ).
\end{enumerate}

\(g\) ) The characteristic polynomial of the automorphism in \(d\) ) (i) is

\[
\mathrm{A} (\mathrm{T}) = (\mathrm{T} - 1) ^ {l} \prod_ {\alpha \in \mathrm{R}} (\mathrm{T} - \zeta^ {\alpha (H)}).
\]

That of the automorphism \(s\) in \(e\) ) is

\[
\mathrm{B} (\mathrm{T}) = \left(\mathrm{T} ^ {h} - 1\right) ^ {l} \prod_ {i = 1} ^ {l} \left(\mathrm{T} - \zeta^ {m _ {i}}\right) \quad \left(m _ {i} \text { being   the   exponents   of } \mathrm{R}\right).
\]

(Use Prop. 33 (iv) of Chap. VI, §1, no. 11.) Deduce from the relation \(\mathrm{A(T)} = \mathrm{B(T)}\) that, for all \(j \geq 1\) , the number of \(i\) such that \(m_i \geq j\) is equal to the number of \(\alpha\in R_{+}\) such that \(\alpha(H)=j\) ; hence recover the result of Exerc. 6 c) of Chap. VI, §4. \(^{12}\)

\begin{enumerate}
\def\labelenumi{\arabic{enumi})}
\setcounter{enumi}{5}
\item
  Assume that \(k = \mathbf{R}\) or \(\mathbf{C}\) . Let \(G\) be the Lie group \(\mathrm{Aut}_0(\mathfrak{g})\) . Show that an element \(a \in G\) is regular (in the sense of Chap. VII, §4, no. 2) if and only if it satisfies the conditions of Exerc. 5 a).
\item
  Assume that g is splittable. Let \(\mathrm{B}(\mathfrak{g})\) be the canonical basis of the canonical Cartan subalgebra of g (no. 3, Remark 2). If \(s \in \mathrm{Aut}(\mathfrak{g})\) , s induces a permutation of \(\mathrm{B}(\mathfrak{g})\) ; denote by \(\operatorname{sgn}(s)\) the sign of this permutation. Show that s operates on \(\bigwedge^{n} g\) (with \(n = \dim g\) ) by
\end{enumerate}

\[
x \mapsto \operatorname{sgn} (s). x.
\]

¶8) Let \(\bar{k}\) be an algebraic closure of k, and \(\bar{g} = \bar{k} \otimes_{k} g\) . The Galois group \(\operatorname{Gal}(\bar{k}/k)\) operates naturally on \(\bar{g}\) , on the canonical Cartan subalgebra of \(\bar{g}\) (no. 3, Remark 2), as well as on its root system \(\bar{R}\) and its canonical basis \(\bar{B}\) . We thus obtain a continuous homomorphism (i.e.~with open kernel)

\[
\pi : \operatorname{Gal} (\bar {k} / k) \to \operatorname{Aut} (\bar {\mathrm{R}}, \bar {\mathrm{B}}).
\]

Show that \(\mathfrak{g}\) is splittable if and only if the following two conditions are satisfied:

\begin{enumerate}
\def\labelenumi{(\roman{enumi})}
\item
  The homomorphism \(\pi\) is trivial.
\item
  g has a Borel subalgebra b.
\end{enumerate}

(Show that a Cartan subalgebra \(\mathfrak{h}\) contained in \(\mathfrak{b}\) is splitting if and only if \(\pi\) is trivial.)

\begin{enumerate}
\def\labelenumi{\arabic{enumi})}
\setcounter{enumi}{8}
\tightlist
\item
  Let R be a reduced root system, B a basis of R, and \((\mathfrak{g}_{0}, \mathfrak{h}_{0}, \mathrm{B}, (X_{\alpha})_{\alpha \in \mathrm{B}})\) a corresponding framed semi-simple Lie algebra (§4). Let \(\bar{k}\) be an algebraic closure of k, and \(\rho : \operatorname{Gal}(\bar{k}/k) \to \operatorname{Aut}(\mathrm{R}, \mathrm{B})\) a continuous homomorphism (cf.~Exerc. 8); if \(\sigma \in \operatorname{Gal}(\bar{k}/k)\) , denote by \(\rho_{\sigma}\) the k-linear automorphism of \(\bar{g}_{0} = \bar{k} \otimes_{k} g_{0}\) such that \(\rho_{\sigma}(X_{\alpha}) = X_{\rho(\sigma)\alpha}\) . On the other hand, the natural operation of \(\operatorname{Gal}(\bar{k}/k)\) on \(\bar{k}\) can be extended to an operation on \(\bar{g}_{0}\) . Let g be the subset of \(\bar{g}_{0}\) consisting of the elements x such that \(\rho_{\sigma}(x) = \sigma^{-1}.x\) for all \(\sigma \in \operatorname{Gal}(\bar{k}/k)\) .
\end{enumerate}

\begin{enumerate}
\def\labelenumi{\alph{enumi})}
\item
  Show that \(\mathfrak{g}\) is a \(k\) -Lie subalgebra of \(\bar{\mathfrak{g}}_0\) , and that the injection of \(\mathfrak{g}\) into \(\bar{\mathfrak{g}}_0\) extends to an isomorphism from \(\bar{k} \otimes_k \mathfrak{g}\) to \(\bar{\mathfrak{g}}_0\) . In particular, \(\mathfrak{g}\) is semi-simple.
\item
  Let \(\mathfrak{b}_0\) be the subalgebra of \(\mathfrak{g}_0\) generated by \(\mathfrak{h}_0\) and the \(X_{\alpha}\) . Put
\end{enumerate}

\[
\bar {\mathfrak {b}} _ {0} = \bar {k} \otimes_ {k} \mathfrak {b} _ {0}, \quad \mathfrak {h} = \mathfrak {g} \cap \bar {\mathfrak {b}} _ {0}, \quad \bar {\mathfrak {h}} _ {0} = \bar {k} \otimes_ {k} \mathfrak {h} _ {0}, \quad \mathfrak {h} = \mathfrak {g} \cap \bar {\mathfrak {h}} _ {0}.
\]

Show that \(\bar{k}\otimes_{k}b=\bar{b}_{0}\) and that \(\bar{k}\otimes_{k}h=\bar{h}_{0}\) , so that b is a Borel subalgebra of g and h is a Cartan subalgebra contained in it.

\begin{enumerate}
\def\labelenumi{\alph{enumi})}
\setcounter{enumi}{2}
\tightlist
\item
  Show that the homomorphism \(\pi\) associated to the algebra g (cf.~Exerc. 8) is equal to \(\rho\) .
\end{enumerate}

¶ 10) Let h be a splitting Cartan subalgebra of g, and \((X_{\alpha})_{\alpha\in\mathbb{R}}\) a Chevalley system in \((\mathfrak{g},\mathfrak{h})\) , cf.~§2, no. 4. If \(\alpha\in R\) , put

\[
\theta_ {\alpha} = e ^ {\mathrm{ad} X _ {\alpha}} e ^ {\mathrm{ad} X _ {- \alpha}} e ^ {\mathrm{ad} X _ {\alpha}},
\]

and denote by \(\overline{W}\) the subgroup of \(\operatorname{Aut}_{e}(\mathfrak{g},\mathfrak{h})\) generated by the \(\theta_{\alpha}\) .

\begin{enumerate}
\def\labelenumi{\alph{enumi})}
\item
  Show that \(\varepsilon (\overline{\mathrm{W}}) = \mathrm{W}(\mathrm{R})\) .
\item
  Let \(s \in \overline{\mathrm{W}}\) , and let \(w = \varepsilon(s)\) . Show that
\end{enumerate}

\[
s (X _ {\alpha}) = \pm X _ {w (\alpha)} \quad \text {   for   all   } \alpha \in \mathrm{R}.
\]

(Use Exerc. 5 of §2.)

\begin{enumerate}
\def\labelenumi{\alph{enumi})}
\setcounter{enumi}{2}
\tightlist
\item
  Let M be the kernel of \(\varepsilon : \overline{\mathrm{W}} \to \mathrm{W}(\mathrm{R})\) . Show that M is contained in the subgroup of \(f(\mathrm{T}_{\mathrm{Q}})\) consisting of the elements \(f(\varphi)\) such that \(\varphi^2 = 1\) . Show that M contains the elements \(f(\varphi_{\alpha})\) defined by \(\varphi_{\alpha}(\beta) = (-1)^{\langle \beta, \alpha^{\vee} \rangle}\) (remark that \(\theta_{\alpha}^{2} = f(\varphi_{\alpha})\) ).
\end{enumerate}

\(d)\) *Let \(\varphi\in\mathrm{Hom}(\mathbf{Q},\{\pm1\})\) . Show that \(f(\varphi)\) belongs to M if and only if \(\varphi\) extends to a homomorphism from P to \(\{\pm1\}\) . (Sufficiency follows from the fact that M contains the \(f(\varphi_{\alpha})\) . To prove necessity, reduce to the case in which \(k=\mathbf{Q}\) , and use the fact that M is contained in \(f(\mathrm{T}_{\mathrm{Q}})\cap\mathrm{Aut}_{e}(\mathfrak{g})=\mathrm{Im}(\mathrm{T}_{\mathrm{P}})\) , cf.~§7, Exerc. 26 \(d)\) .) Deduce that M is isomorphic to the dual of the group \(\mathrm{Q}/(\mathrm{Q}\cap2\mathrm{P}).^{13}\)

\begin{enumerate}
\def\labelenumi{\arabic{enumi})}
\setcounter{enumi}{10}
\tightlist
\item
  With the notations of no. 2, assume that \(k\) is non-discrete and locally compact, hence isomorphic to \(\mathbf{R}\) , \(\mathbf{C}\) or a finite extension of \(\mathbf{Q}_p\) (Commutative Algebra, Chap. VI, §9, no. 3). For all \(n \geq 1\) , the quotient \(k^* / k^{*n}\) is finite (cf.~Commutative Algebra, Chap. VI, §9, Exerc. 3 for the ultrametric case). Deduce that the quotients \(T_Q / \text{Im}(T_P)\) and \(\text{Aut}(\mathfrak{g}) / \text{Aut}_e(\mathfrak{g})\) are finite.
\end{enumerate}

When \(k = \mathbf{R}\) , show that \(\mathrm{T_Q / Im(T_P)}\) is isomorphic to the dual of the \(\mathbf{F}_{2}\) -vector space \((\mathrm{Q} \cap 2\mathrm{P}) / 2\mathrm{Q}\) . When \(k = \mathbf{C}\) , we have \(\mathrm{T_Q = Im(T_P)}\) ; this is the integral subgroup of the Lie group \(\operatorname{Aut}(\mathfrak{g})\) with Lie algebra \(\mathfrak{h}\) .

¶ 12) Let \(\mathfrak{h}\) be a splitting Cartan subalgebra of \(\mathfrak{g}\) , A a subset of \(\mathfrak{h}\) , and \(s \in \operatorname{Aut}(\mathfrak{g})\) such that \(s\mathrm{A} = \mathrm{A}\) . Show that there exists \(t \in \operatorname{Aut}(\mathfrak{g}, \mathfrak{h})\) such that \(t|\mathrm{A} = s|\mathrm{A}\) and \(ts^{-1} \in \operatorname{Aut}_e(\mathfrak{g})\) . (Let \(\mathfrak{a}\) be the commutant of A in \(\mathfrak{g}\) ; this is a reductive subalgebra of \(\mathfrak{g}\) , of which \(sh\) and \(\mathfrak{h}\) are splitting Cartan subalgebras; deduce the existence of \(u \in \operatorname{Aut}_e(\mathfrak{a})\) such that \(ush = \mathfrak{h}\) ; show that there exists \(v \in \operatorname{Aut}_e(\mathfrak{g})\) extending \(u\) such that \(v|\mathrm{A} = \mathrm{Id}_{\mathrm{A}}\) ; take \(t = vs.\) ) Deduce that, if \(s \in \operatorname{Aut}_0(\mathfrak{g})\) , there exists \(w \in \mathrm{W}(\mathrm{R})\) such that \(w|\mathrm{A} = s|\mathrm{A}\) .

¶13) Let \((\mathfrak{g},\mathfrak{h},\mathrm{B},(X_{\alpha})_{\alpha\in\mathrm{B}})\) be a framed semi-simple Lie algebra, R the root system of \((\mathfrak{g},\mathfrak{h})\) , \(\Delta\) the corresponding Dynkin graph, and \(\varPhi\) a subgroup of

\[
\operatorname{Aut} (\mathrm{R}, \mathrm{B}) = \operatorname{Aut} (\Delta).
\]

If \(s \in \varPhi\) , extend \(s\) to an automorphism of \(\mathfrak{g}\) by the conditions

\[
s (X _ {\alpha}) = X _ {s \alpha}, \quad \text { and } \quad s (H _ {\alpha}) = H _ {s \alpha} \quad \text { for   all } \alpha \in \mathrm{B}, \text { cf.   Prop. } 1.
\]

We thus identify \(\varPhi\) with a subgroup of \(\operatorname{Aut}(\mathfrak{g},\mathfrak{h})\) ; denote by \(\tilde{\mathfrak{g}}\) (resp. \(\tilde{\mathfrak{h}}\) ) the subalgebra of g (resp. h) consisting of the elements invariant under \(\varPhi\) .

\begin{enumerate}
\def\labelenumi{\alph{enumi})}
\tightlist
\item
  Let \(\alpha \in \mathbf{B}\) , and let \(\mathbf{X} = \varPhi.\alpha\) . Show, by using the Plates in Chap. VI, that only two cases are possible:
\end{enumerate}

\begin{enumerate}
\def\labelenumi{(\roman{enumi})}
\item
  every element of \(X\) distinct from \(\alpha\) is orthogonal to \(\alpha\) ;
\item
  there exists a unique element \(\alpha'\) of \(X - \{\alpha\}\) that is not orthogonal to \(\alpha\) , and \(n(\alpha, \alpha') = n(\alpha', \alpha) = -1\) .
\end{enumerate}

\begin{enumerate}
\def\labelenumi{\alph{enumi})}
\setcounter{enumi}{1}
\tightlist
\item
  Let i be the restriction map \(h^{*} \to \tilde{h}^{*}\) , and let \(\tilde{\mathbf{B}} = i(\mathbf{B})\) ; the map \(\mathbf{B} \to \tilde{\mathbf{B}}\) identifies \(\tilde{B}\) with \(B/\varPhi\) . Show that \(\tilde{g}\) is a semi-simple Lie algebra, that \(\tilde{h}\) is a splitting Cartan subalgebra of it, and that \(\tilde{B}\) is a basis of \(\mathrm{R}(\tilde{\mathfrak{g}}, \tilde{\mathfrak{h}})\) . (Observe that \(\tilde{B}\) is contained in \(\mathrm{R}(\tilde{\mathfrak{g}}, \tilde{\mathfrak{h}})\) , and that every element of \(\mathrm{R}(\tilde{\mathfrak{g}}, \tilde{\mathfrak{h}})\) is a linear combination of elements of \(\tilde{B}\) with integer coefficients of the same sign.) If \(\tilde{\alpha} \in \tilde{B}\) , the corresponding inverse root \(H_{\tilde{\alpha}} \in \tilde{h}\) is given by
\end{enumerate}

\[
\begin{array}{l} H _ {\tilde {\alpha}} = \sum_ {i (\alpha) = \tilde {\alpha}} H _ {\alpha} \quad \text {in case (i) of a)} \\ H _ {\tilde {\alpha}} = 2 \sum_ {i (\alpha) = \tilde {\alpha}} H _ {\alpha} \quad \text {in case (ii) of a)}, \end{array}
\]

where the summation is over those elements \(\alpha \in \mathrm{B}\) such that \(i(\alpha) = \tilde{\alpha}\) . If \(\beta \in \mathrm{B}\) has image \(\tilde{\beta} \in \tilde{\mathrm{B}}\) , then

\[
  n (\tilde {\beta}, \tilde {\alpha}) = \sum_ {i (\alpha) = \tilde {\alpha}} n (\beta , \alpha) \quad \text { in   case   (i) }
\]

\[
n (\tilde {\beta}, \tilde {\alpha}) = 2 \sum_ {i (\alpha) = \tilde {\alpha}} n (\beta , \alpha) \quad \text { in   case   (ii) }.
\]

Deduce how the Dynkin graph of \(\mathrm{R}(\tilde{\mathfrak{g}},\tilde{\mathfrak{h}})\) is determined starting from the pair \((\varDelta,\varPhi)\) .

\begin{enumerate}
\def\labelenumi{\alph{enumi})}
\setcounter{enumi}{2}
\tightlist
\item
  Show that, if \(\mathfrak{g}\) is simple, so is \(\tilde{\mathfrak{g}}\) .
\end{enumerate}

If \(\mathfrak{g}\) is of type \(A_l\) , \(l \geq 2\) , and \(\varPhi\) is of order 2, then \(\tilde{\mathfrak{g}}\) is of type \(B_{l/2}\) if \(l\) is even, and of type \(C_{(l+1)/2}\) if \(l\) is odd.

If \(\mathfrak{g}\) is of type \(D_l\) , \(l \geq 4\) , and \(\varPhi\) is of order 2, then \(\tilde{\mathfrak{g}}\) is of type \(B_{l-1}\) .

If \(\mathfrak{g}\) is of type \(\mathrm{D}_4\) , and \(\varPhi\) is of order 3 or 6, then \(\tilde{\mathfrak{g}}\) is of type \(\mathrm{G}_2\) .

If g is of type \(E_{6}\) , and \(\Phi\) is of order 2, then \(\tilde{g}\) is of type \(F_{4}\) .

\exercisesection{\S~6}

\begin{enumerate}
\def\labelenumi{\arabic{enumi})}
\item
  Show that \(Z(\lambda)\) can be defined as a quotient of the representation \(\rho_{\lambda}\) of §4, Exerc. 1.
\item
  Let \(\mu\) be a weight of \(Z(\lambda)\) (resp. \(E(\lambda)\) ). Show that there exists a sequence of weights \(\mu_0, \ldots, \mu_n\) of \(Z(\lambda)\) (resp. of \(E(\lambda)\) ) such that \(\mu_0 = \lambda, \mu_n = \mu\) , and \(\mu_{i-1} - \mu_i \in B\) for \(1 \leq i \leq n\) .
\item
  Assume that g is simple, and denote by \(\tilde{\alpha}\) the highest root of R. The module \(\mathrm{E}(\tilde{\alpha})\) is isomorphic to g, equipped with the adjoint representation. If C is the Casimir element associated to the Killing form of g, the image of C in \(\mathrm{End}\mathrm{E}(\tilde{\alpha})\) is the identity (cf.~Chap. I, §3, no. 7, Prop. 12). Deduce, by using the Cor. of Prop. 7, that
\end{enumerate}

\[
\varPhi_ {\mathrm{R}} (\tilde {\alpha}, \tilde {\alpha} + 2 \rho) = 1,
\]

where \(\Phi_{\mathrm{R}}\) is the canonical bilinear form on \(\mathfrak{h}^*\) (Chap. VI, §1, no. 12).

\begin{enumerate}
\def\labelenumi{\arabic{enumi})}
\setcounter{enumi}{3}
\tightlist
\item
  We use the notations of no. 4.
\end{enumerate}

\begin{enumerate}
\def\labelenumi{\alph{enumi})}
\item
  Let \(m \in \mathbf{N}\) . Give \(\mathbf{N}^m\) the product order. Show that, for any subset \(S\) of \(\mathbf{N}^m\) , the set of minimal elements of \(S\) is finite.
\item
  Let \(\alpha_{1},\ldots ,\alpha_{m}\) be distinct elements of R, \(X_{i}\in \mathfrak{g}^{\alpha_{i}} - \{0\}\) , S the set of non-zero sequences \((p_i)\in \mathbf{N}^m\) such that \(\sum p_i\alpha_i = 0\) , M the set of minimal elements of S. Then \(\mathfrak{h}\) and the \(X_1^{p_1}\dots X_m^{p_m}\) , where \((p_1,\dots ,p_m)\in \mathrm{M}\) , generate the algebra \(\mathrm{U}^0\) .
\item
  Show that \(U^{0}\) is both a left- and right-noetherian algebra. (Give \(U^{0}\) the filtration induced by that of \(\mathrm{U}(\mathfrak{g})\) , and show, by using a) and b), that gr \(U^{0}\) is commutative of finite type.)
\end{enumerate}

\(d\) ) Show that, for all \(\lambda \in \mathfrak{h}^*\) , \(\mathrm{U}^{\lambda}\) is a left (resp. right) \(\mathrm{U}^0\) -module of finite type.

\begin{enumerate}
\def\labelenumi{\alph{enumi})}
\setcounter{enumi}{4}
\tightlist
\item
  Let V be a simple g-module such that \(V = \bigoplus_{\lambda \in b^{*}} V_{\lambda}\) . If one of the \(V_{\lambda} \neq 0\) is finite dimensional, then all of the \(V_{\lambda}\) are finite dimensional. (Use d).
\end{enumerate}

\begin{enumerate}
\def\labelenumi{\arabic{enumi})}
\setcounter{enumi}{4}
\tightlist
  \item
  Show that, if \(\mathfrak{g} = \mathfrak{sl}(2,k)\) , the modules \(Z(\lambda)\) of this paragraph are isomorphic to the modules \(Z(\lambda)\) of §1, Exerc. 2.
\end{enumerate}

\exercisesection{\S~7}

All the g-modules considered (except those in Exerc. 14 and 15) are assumed to be finite dimensional.

\begin{enumerate}
\def\labelenumi{\arabic{enumi})}
\tightlist
\item
  Let \(\omega \in \mathrm{P}_{++}\) ; denote by \(\mathrm{S}(\omega)\) the set of weights of \(\mathrm{E}(\omega)\) , in other words the smallest R-saturated subset of P containing \(\omega\) (Prop. 5). If \(\lambda \in \mathrm{S}(\omega)\) , we have \(\lambda \equiv \omega\) (mod. Q). Conversely, let \(\lambda \in \mathrm{P}\) be such that \(\lambda \equiv \omega\) (mod. Q); prove the equivalence of the following properties:
\end{enumerate}

\begin{enumerate}
\def\labelenumi{(\roman{enumi})}
\item
  \(\lambda \in S(\omega);\)
\item
  \(\omega - w\lambda \in Q_{+}\) for all \(w \in W\) ;
\item
  \(\lambda\) belongs to the convex hull of W. \(\omega\) in \(\mathfrak{h}_{\mathbf{R}}^{*}\) .
\end{enumerate}

(To prove that (iii) \(\Longrightarrow\) (ii), remark that \(\omega - w\omega\) is a linear combination of elements of \(\mathrm{R}_{+}\) with coefficients \(\geq 0\) ; deduce, by convexity, that \(\omega - w\lambda\) has the same property; since \(\omega - w\lambda\) belongs to Q, this implies that \(\omega - w\lambda \in \mathrm{Q}_{+}\) . To prove that (ii) \(\Longrightarrow\) (i), choose \(w\) such that \(w\lambda \in \mathrm{P}_{++}\) , and apply Cor. 2 of Prop. 3. The implication (i) \(\Longrightarrow\) (iii) is immediate.)

\begin{enumerate}
\def\labelenumi{\arabic{enumi})}
\setcounter{enumi}{1}
\tightlist
\item
  Let \((\mathrm{R}_{i})_{i\in\mathrm{I}}\) be the family of irreducible components of R, and \(g=\prod_{i\in I}g_{i}\) the corresponding decomposition of g into a product of simple algebras. Identify P with the product of the \(\mathrm{P}(\mathrm{R}_{i})\) , and give each \(\mathrm{P}(\mathrm{R}_{i})\) the order relation defined by the basis \(B_{i}=B\cap R_{i}\) .
\end{enumerate}

\begin{enumerate}
\def\labelenumi{\alph{enumi})}
\item
  Let \(\omega = (\omega_{i})_{i\in \mathrm{I}}\) be an element of \(\mathrm{P}_{++} = \prod_{i\in \mathrm{I}}\mathrm{P}_{++}(\mathrm{R}_i)\) . Show that the simple \(\mathfrak{g}\) -module \(\operatorname {E}(\omega)\) is isomorphic to the tensor product of the simple \(\mathfrak{g}_i\) -modules \(\operatorname {E}(\omega_i)\) .
\item
  Let \(\mathcal{M}\) (resp. \(M_{i}\) ) be the set of elements of \(P_{++}\) (resp. of \(P_{++}(R_{i})\) ) having the equivalent properties (i), (ii), (iii), (iv) of Prop. 6 and 7. Show that \(M = \prod_{i \in I} M_{i}\) , in other words that \(\omega \in M\) if and only if, for all \(i \in I\) , \(\omega_{i}\) is either zero or a minuscule weight of \(R_{i}\) . Deduce that M is a system of representatives in P of the elements of P/Q.
\item
  Let E be a simple g-module, and X its set of weights. Show that X contains a unique element of M, and that the multiplicity of this element is equal to the upper bound of the multiplicities of the elements of X.
\end{enumerate}

\begin{enumerate}
\def\labelenumi{\arabic{enumi})}
\setcounter{enumi}{2}
\tightlist
\item
  \begin{enumerate}
  \def\labelenumii{\alph{enumii})}
  \tightlist
  \item
    Let E be a g-module. Show the equivalence of the conditions:
  \end{enumerate}
\end{enumerate}

\begin{enumerate}
\def\labelenumi{(\roman{enumi})}
\item
  The rank of the semi-direct product of \(\mathfrak{g}\) by E is strictly larger than that of \(\mathfrak{g}\) .
\item
  0 is a weight of E.
\item
  There exists a weight of \(\mathrm{E}\) that is radical (i.e.~belongs to \(\mathrm{Q}\) ).
\end{enumerate}

\begin{enumerate}
\def\labelenumi{\alph{enumi})}
\setcounter{enumi}{1}
\tightlist
\item
  Assume that \(\mathrm{E}\) is simple. Show that (i), (ii), (iii) are equivalent to
\end{enumerate}

\begin{enumerate}
\def\labelenumi{(\roman{enumi})}
\setcounter{enumi}{3}
\tightlist
\item
  The highest weight of E is radical.
\end{enumerate}

If these conditions are satisfied, there exists no non-zero invariant alternating bilinear form. (Use Prop. 12 and Prop. 1 (ii) of §6.)

\begin{enumerate}
\def\labelenumi{\arabic{enumi})}
\setcounter{enumi}{3}
\item
  Let \(k'\) be an extension of \(k\) , and \(\mathfrak{g}' = \mathfrak{g}_{(k')}\) . Show that every \(\mathfrak{g}'\) -module arises, by extension of scalars, from a \(\mathfrak{g}\) -module that is unique up to isomorphism.
\item
  Let \(\mathbf{E}\) be a \(\mathfrak{g}\) -module. Show the equivalence of the conditions:
\end{enumerate}

\begin{enumerate}
\def\labelenumi{(\roman{enumi})}
\item
  E is faithful (i.e.~the canonical map from g to gl(E) is injective).
\item
  Every root of g is the difference between two weights of E.
\end{enumerate}

¶6) Let \(\varphi\) be an involutive automorphism of \(\mathfrak{g}\) whose restriction to \(\mathfrak{h}\) is -Id. Show that, if E is a \(\mathfrak{g}\) -module, there exists a non-degenerate symmetric bilinear form \(\varPsi\) on E such that

\[
\varPsi (x. a, b) + \varPsi (a, \varphi (x). b) = 0 \quad \text {   for   } x \in \mathfrak {g}, a, b \in \mathrm{E}.
\]

(Reduce to the case in which E is simple. Show that the transform of E by \(\varphi\) is isomorphic to the dual \(E^{*}\) of E, and deduce the existence of a non-degenerate bilinear form \(\Psi\) satisfying the conditions above. Show next that, if e is a primitive vector in E, then \(\Psi(e,e)\neq0\) . Deduce that \(\Psi\) is symmetric.)

\begin{enumerate}
\def\labelenumi{\arabic{enumi})}
\setcounter{enumi}{6}
\tightlist
\item
  If \(\lambda \in \mathrm{P}_{++}\) , denote by \(\rho_{\lambda}:\mathrm{U}(\mathfrak{g})\to \mathrm{End}(\mathrm{E}(\lambda))\) the representation defined by the simple module \(\mathrm{E}(\lambda)\) . We have \(\mathrm{Im}(\rho_{\lambda}) = \mathrm{End}(\mathrm{E}(\lambda));\) put \(\mathfrak{m}_{\lambda} = \mathrm{Ker}(\rho_{\lambda})\) . \(a)\) Show that the \(\mathfrak{m}_{\lambda}\) are pairwise distinct, and that they are the only two-sided ideals \(\mathfrak{m}\) of \(\mathrm{U}(\mathfrak{g})\) such that \(\mathrm{U}(\mathfrak{g}) / \mathfrak{m}\) is a finite dimensional simple \(k\) -algebra.
\end{enumerate}

\begin{enumerate}
\def\labelenumi{\alph{enumi})}
\setcounter{enumi}{1}
\item
  If I is a finite subset of \(P_{++}\) , put \(m_{I} = \bigcap_{\lambda \in I} m_{\lambda}\) . Show that the canonical map \(U(\mathfrak{g}) / \mathfrak{m}_{I} \to \prod_{\lambda \in I} U(\mathfrak{g}) / \mathfrak{m}_{\lambda}\) is an isomorphism, and that every two-sided ideal of \(U(\mathfrak{g})\) of finite codimension is equal to exactly one of the \(m_{I}\) .
\item
  Show that the principal anti-automorphism of \(\mathrm{U}(\mathfrak{g})\) transforms \(\mathfrak{m}_{\lambda}\) to \(\mathfrak{m}_{\lambda^*}\) , where \(\lambda^{*} = -w_{0}\lambda\) (cf.~Prop. 11).
\end{enumerate}

¶8) Let g be a semi-simple Lie algebra; exceptionally, we do not assume in this exercise that g is split. Let \(\bar{k}\) be an algebraic closure of k and let

\[
\pi : \operatorname{Gal} (\bar {k} / k) \to \operatorname{Aut} (\bar {\mathrm{R}}, \bar {\mathrm{B}})
\]

be the homomorphism defined in §5, Exerc. 8. Let \(\operatorname{Gal}(\bar{k}/k)\) operate, via \(\pi\) , on the set \(\bar{P}_{++}\) of dominant weights of \(\bar{R}\) relative to \(\bar{B}\) ; let \(\Omega\) be a system of representatives of the elements of the quotient \(\bar{P}_{++}/\operatorname{Gal}(\bar{k}/k)\) .

\begin{enumerate}
\def\labelenumi{\alph{enumi})}
\item
  Put \(\bar{g} = \bar{k} \otimes_{k} g\) . If I is a finite subset of \(\bar{P}_{++}\) stable under \(\operatorname{Gal}(\bar{k}/k)\) , the two-sided ideal \(\bar{m}_{I}\) of \(\operatorname{U}(\bar{g})\) associated to I (cf.~Exerc. 7) is of the form \(\bar{k} \otimes_{k} m_{I}\) , where \(m_{I}\) is a two-sided ideal of \(\operatorname{U}(g)\) . Show that every two-sided ideal of \(\operatorname{U}(g)\) of finite codimension is obtained in this way exactly once.
\item
  Let \(\omega \in \Omega\) , \(\mathrm{I}(\omega)\) its orbit under \(\operatorname{Gal}(\bar{k}/k)\) , and \(\mathrm{G}_{\omega}\) the stabilizer of \(\omega\) ; let \(k_{\omega}\) be the sub-extension of \(\bar{k}\) corresponding to \(\mathrm{G}_{\omega}\) by Galois theory. Show that \(\mathrm{U}(\mathfrak{g})/\mathfrak{m}_{\mathrm{I}(\omega)}\) is a simple algebra whose centre is isomorphic to \(k_{\omega}\) . Every two-sided ideal \(\mathfrak{m}\) of \(\mathrm{U}(\mathfrak{g})\) such that \(\mathrm{U}(\mathfrak{g})/\mathfrak{m}\) is a finite dimensional simple algebra is equal to exactly one of the \(\mathfrak{m}_{\mathrm{I}(\omega)}\) .
\item
  The group \(\operatorname{Gal}(\bar{k}/k)\) operates, via \(\pi\) , on the ring \(\operatorname{R}(\mathfrak{g})\) ; denote by \(\operatorname{R}(\mathfrak{g})^{\mathrm{inv}}\) the subring of \(\operatorname{R}(\mathfrak{g})\) consisting of the elements invariant under \(\operatorname{Gal}(\bar{k}/k)\) . Show that the map \([E] \to [\bar{k} \otimes_{k} E]\) extends to an injective homomorphism from \(\operatorname{R}(\mathfrak{g})\) to \(\operatorname{R}(\mathfrak{g})^{\mathrm{inv}}\) whose cokernel is a torsion group; this is an isomorphism if and only if, for all \(\omega \in \Omega\) , \(\operatorname{U}(\mathfrak{g})/\mathfrak{m}_{\operatorname{I}(\omega)}\) is an algebra of matrices over \(k_{\omega}\) ; show that this is the case when g has a Borel subalgebra. \(^{14}\)
\end{enumerate}

\begin{enumerate}
\def\labelenumi{\arabic{enumi})}
\setcounter{enumi}{8}
\tightlist
\item
  Let \(\mathfrak{a}_1\) and \(\mathfrak{a}_2\) be Lie algebras. Show that there exists a unique homomorphism
\end{enumerate}

\[
f: \mathrm{R} (\mathfrak {a} _ {1}) \otimes_ {\mathbf {Z}} \mathrm{R} (\mathfrak {a} _ {2}) \to \mathrm{R} (\mathfrak {a} _ {1} \times \mathfrak {a} _ {2})
\]

such that \(f([\mathrm{E}_1] \otimes [\mathrm{E}_2]) = [\mathrm{E}_1 \otimes \mathrm{E}_2]\) if \(\mathrm{E}_i\) is an \(\mathfrak{a}_i\) -module ( \(i = 1, 2\) ). Show that \(f\) is injective, and that it is bijective if \(\mathfrak{a}_1\) and \(\mathfrak{a}_2\) are splittable semi-simple.

\begin{enumerate}
\def\labelenumi{\arabic{enumi})}
\setcounter{enumi}{9}
\tightlist
\item
  Let \(\Gamma\) be a subgroup of \(P\) containing \(Q\) . Such a subgroup is stable under \(W\) .
\end{enumerate}

\begin{enumerate}
\def\labelenumi{\alph{enumi})}
\item
  Show that, if \(\lambda \in \mathrm{P}_{++} \cap \Gamma\) , every weight of \(\operatorname{E}(\lambda)\) belongs to \(\Gamma\) .
\item
  Let \(\mathrm{R}_{\Gamma}(\mathfrak{g})\) be the subgroup of \(\mathrm{R}(\mathfrak{g})\) with basis the \([\lambda]\) , with \(\lambda \in \mathrm{P}_{++} \cap \Gamma\) . If \(\mathrm{E}\) is a \(\mathfrak{g}\) -module, \([\mathrm{E}] \in \mathrm{R}_{\Gamma}(\mathfrak{g})\) if and only if the weights of \(\mathrm{E}\) belong to \(\Gamma\) . Deduce that \(\mathrm{R}_{\Gamma}(\mathfrak{g})\) is a subring of \(\mathrm{R}(\mathfrak{g})\) .
\item
  Show that the homomorphism \(\operatorname{ch} : \mathrm{R}_{\Gamma}(\mathfrak{g}) \to \mathbf{Z}[\Gamma]\) is an isomorphism from \(\mathrm{R}_{\Gamma}(\mathfrak{g})\) to the subring of \(\mathbf{Z}[\Gamma]\) consisting of the elements invariant under W. (Use Th. 2 (ii).)
\end{enumerate}

\(d\) ) Describe \(\mathrm{R}(\mathfrak{g})\) and \(\mathrm{R}_{\Gamma}(\mathfrak{g})\) explicitly when \(\mathfrak{g} = \mathfrak{sl}(2,k)\) and \(\Gamma = \mathrm{Q}\) .

¶11) The notations are those of no. 7. For any integer \(m \geq 1\) , denote by \(\Psi^{m}\) the endomorphism of \(\mathbf{Z}[\Delta]\) that takes \(e^{\lambda}\) to \(e^{m\lambda}\) . We have \(\Psi^{1} = \mathrm{Id}\) and \(\Psi^{m} \circ \Psi^{n} = \Psi^{mn}\) .

Let E be a finite dimensional \(\Delta\) -graded vector space. For all \(n \geq 0\) , denote by \(a_{n}E\) (resp. \(s_{n}E\) ) the nth exterior (resp. symmetric) power of E, equipped with its natural grading.

\begin{enumerate}
\def\labelenumi{\alph{enumi})}
\tightlist
\item
  Show that
\end{enumerate}

\[
n \operatorname{ch} (s _ {n} \mathrm{E}) = \sum_ {m = 1} ^ {n} \Psi^ {m} (\operatorname{ch} (\mathrm{E})) \operatorname{ch} (s _ {n - m} \mathrm{E})
\]

and

\[
n \operatorname{ch} (a _ {n} \mathrm{E}) = \sum_ {m = 1} ^ {n} (- 1) ^ {m - 1} \varPsi^ {m} (\operatorname{ch} (\mathrm{E})) \operatorname{ch} (a _ {n - m} \mathrm{E}).
\]

Deduce that \(\operatorname{ch}(s_{n}\operatorname{E})\) and \(\operatorname{ch}(a_{n}\operatorname{E})\) can be expressed as polynomials, with rational coefficients, in the \(\varPsi^{m}(\operatorname{ch}(\operatorname{E}))\) , \(1 \leq m \leq n\) . For example:

\[
\mathrm{ch} (s _ {2} \mathrm{E}) = \frac {1}{2} \mathrm{ch} (\mathrm{E}) ^ {2} + \frac {1}{2} \varPsi^ {2} (\mathrm{ch} (\mathrm{E}))
\]

\[
\mathrm{ch} (a _ {2} \mathrm{E}) = \frac {1}{2} \mathrm{ch} (\mathrm{E}) ^ {2} - \frac {1}{2} \varPsi^ {2} (\mathrm{ch} (\mathrm{E})).
\]

\begin{enumerate}
\def\labelenumi{\alph{enumi})}
\setcounter{enumi}{1}
\tightlist
\item
  Prove the following identities (in the algebra of formal power series in a variable T, with coefficients in \(Q[\Delta]\) ):
\end{enumerate}

\[
\sum_ {n = 0} ^ {\infty} \mathrm{ch} (s _ {n} \mathrm{E}) \mathrm{T} ^ {n} = \exp \left\{\sum_ {m = 1} ^ {\infty} \varPsi^ {m} (\mathrm{ch(E)}) \mathrm{T} ^ {m} / m \right\}
\]

and

\[
\sum_ {n = 0} ^ {\infty} \mathrm{ch} (a _ {n} \mathrm{E}) \mathrm{T} ^ {n} = \exp \left\{\sum_ {m = 1} ^ {\infty} (- 1) ^ {m - 1} \varPsi^ {m} (\mathrm{ch(E)}) \mathrm{T} ^ {m} / m \right\}.
\]

\begin{enumerate}
\def\labelenumi{\alph{enumi})}
\setcounter{enumi}{2}
\tightlist
\item
  Assume that \(\Delta\) is a group, so that \(\varPsi^{m}\) can be defined for all \(m\in \mathbf{Z}\) . Show that, if \(\mathrm{E}^*\) is the graded dual of \(\mathrm{E}\) ,
\end{enumerate}

\[
\operatorname{ch} \left(\mathrm{E} ^ {*}\right) = \Psi^ {- 1} (\operatorname{ch} (\mathrm{E})).
\]

\begin{enumerate}
\def\labelenumi{\alph{enumi})}
\setcounter{enumi}{3}
\tightlist
\item
  Identify \(\mathrm{R}(\mathfrak{g})\) , by means of ch, with a subring of Z[P]. Show that \(\mathrm{R}(\mathfrak{g})\) is stable under the \(\varPsi^{m}\) , \(m \in Z\) , and that so is the subring \(\mathrm{R}_{\varGamma}(\mathfrak{g})\) defined in the preceding exercise.
\end{enumerate}

\begin{enumerate}
\def\labelenumi{\arabic{enumi})}
\setcounter{enumi}{11}
\tightlist
\item
  Let \(\lambda \in \mathrm{P}_{++}\) and let \(\mathcal{X}_{\lambda}\) be the set of weights of \(\operatorname{E}(\lambda)\) . Show that
\end{enumerate}

\[
\mathcal {X} _ {\lambda} \subset \lambda - \mathrm{P} _ {+ +}
\]

does not hold in general. (Consider, for example, the adjoint representation of \(\mathfrak{sl}(3,k)\) .)

\begin{enumerate}
\def\labelenumi{\arabic{enumi})}
\setcounter{enumi}{12}
\tightlist
\item
  Let \(\lambda = \sum_{\alpha \in B} a_{\alpha} \alpha\) be an element of \(P_{++}\) . For \(n = 0, 1, \ldots\) , denote by \(\mathcal{X}_n\) the set of weights \(\mu\) of \(E(\lambda)\) such that \(\lambda - \mu\) is the sum of \(n\) elements of \(B\) . Let \(s_n\) be the sum of the multiplicities of the elements of \(\mathcal{X}_n\) (as weights of \(E(\lambda)\) ). Let \(T = 2 \sum_{\alpha \in B} a_{\alpha}\) . Show that:
\end{enumerate}

\begin{enumerate}
\def\labelenumi{\alph{enumi})}
\item
  T is an integer \(\geq 0\) .
\item
  \(s_{n}=0\) for n > T, and \(s_{T-n}=s_{n}\) .
\item
  If r is the integer part of T/2, then \(s_{1} \leq s_{2} \leq \cdots \leq s_{r+1}\) .
\end{enumerate}

¶14) Let \(\lambda \in \mathrm{P}_{++}\) , \(\mathrm{F}_{\lambda}\) be the largest proper submodule of \(Z(\lambda)\) and \(v\) a primitive element of \(Z(\lambda)\) of weight \(\lambda\) , cf.~§6, no. 3. Show that

\[
\mathrm{F} _ {\lambda} = \sum_ {\alpha \in \mathrm{B}} \mathrm{U} (\mathfrak {g}) X _ {- \alpha} ^ {\lambda (H _ {\alpha}) + 1} v = \sum_ {\alpha \in \mathrm{B}} \mathrm{U} (\mathfrak {n} _ {-}) X _ {- \alpha} ^ {\lambda (H _ {\alpha}) + 1} v.
\]

\begin{enumerate}
\def\labelenumi{\arabic{enumi})}
\setcounter{enumi}{14}
\tightlist
\item
  Let \(\lambda \in \mathfrak{h}^*\) and let \(v\) (resp. \(v'\) ) be a primitive element of \(\mathrm{Z}(\lambda)\) (resp. of \(\mathrm{E}(\lambda)\) ) of weight \(\lambda\) . Let \(\mathrm{I}\) (resp. \(\mathrm{I}'\) ) be the annihilator of \(v\) (resp. \(v'\) ) in \(\mathrm{U}(\mathfrak{g})\) .
\end{enumerate}

\[
a) \mathrm{I} = \mathrm{U} (\mathfrak {g}) \mathfrak {n} _ {+} + \sum_ {h \in \mathfrak {h}} \mathrm{U} (\mathfrak {g}) (h - \lambda (h)).
\]

\begin{enumerate}
\def\labelenumi{\alph{enumi})}
\setcounter{enumi}{1}
\item
  I' is the largest left ideal of U(g) distinct from U(g) and containing I.
\item
  If \(\lambda \in \mathrm{P}_{++}\) , then
\end{enumerate}

\[
\mathrm{I} ^ {\prime} = \mathrm{I} + \sum_ {\alpha \in \mathrm{B}} \mathrm{U} (\mathfrak {g}) X _ {- \alpha} ^ {\lambda (H _ {\alpha}) + 1} = \mathrm{I} + \sum_ {\alpha \in \mathrm{B}} \mathrm{U} (\mathfrak {n} _ {-}) X _ {- \alpha} ^ {\lambda (H _ {\alpha}) + 1}.
\]

(Use the preceding exercise.)

¶16) Let V and V' be two g-modules. Then V' is said to be subordinate to V if there exists a linear map \(f : V \to V'\) such that:

\(\alpha)\) \(f\) is surjective; \(\beta)\) the image under \(f\) of a primitive element of \(\mathbf{V}\) is either 0 or a primitive element of \(\mathrm{V}^{\prime}\) ; \(\gamma)\) \(f\) is an \(\mathfrak{n}_{-}\) -homomorphism.

\begin{enumerate}
\def\labelenumi{\alph{enumi})}
\item
  Let \(f: V \to V'\) satisfy \(\alpha\) ), \(\beta\) ), \(\gamma\) ). Let v be a primitive element of V. Then the image under f of the g-submodule generated by v is the g-submodule generated by \(f(v)\) .
\item
  Let \(f: \mathrm{V} \to \mathrm{V}'\) satisfy \(\alpha), \beta), \gamma)\) . Let \(\mathrm{W}\) be a \(\mathfrak{g}\) -submodule of \(\mathrm{V}\) . Then \(f(\mathrm{W})\) is a \(\mathfrak{g}\) -submodule of \(\mathrm{V}'\) that is subordinate to \(\mathrm{W}\) .
  \item
  Let \(V = E_{1} \oplus \cdots \oplus E_{s}\) and \(V' = E_{1}' \oplus \cdots E_{s'}\) be decompositions of V and \(V'\) into sums of simple modules. Then \(V'\) is subordinate to V if and only if \(s' \leq s\) and there exists \(\sigma \in S_{s}\) such that \(E_{i}'\) is subordinate to \(\mathrm{E}_{\sigma(i)}\) for \(i = 1, \ldots, s'\) .
\item
  If \(V'\) is subordinate to V and if V is simple, then \(V'\) is simple or reduced to 0.
\item
  Assume that V and \(V'\) are simple. Let \(\lambda\) and \(\lambda'\) be their highest weights. Then \(V'\) is subordinate to V if and only if \(\lambda'(H_{\alpha}) \leq \lambda(H_{\alpha})\) for all \(\alpha \in B\) . (For the sufficiency, use the preceding exercise.)
\end{enumerate}

\begin{enumerate}
\def\labelenumi{\arabic{enumi})}
\setcounter{enumi}{16}
\tightlist
\item
  Let \(\lambda, \mu \in \mathrm{P}_{++}\) and \(\alpha \in \mathrm{B}\) be such that \(\lambda(H_{\alpha}) \geq 1\) and \(\mu(H_{\alpha}) \geq 1\) . Let \(\mathrm{F} = \mathrm{E}(\lambda) \otimes \mathrm{E}(\mu)\) .
\end{enumerate}

\begin{enumerate}
\def\labelenumi{\alph{enumi})}
\item
  Show that \(\dim \mathrm{F}^{\lambda +\mu} = 1\) and \(\dim \mathrm{F}^{\lambda +\mu -\alpha} = 2\) .
\item
  Show that \(X_{\alpha}: F^{\lambda+\mu-\alpha} \to F^{\lambda+\mu}\) is surjective, and that the non-zero elements of its kernel are primitive (remark that, if \(\beta \in B\) is distinct from \(\alpha\) , \(\lambda + \mu - \alpha + \beta\) is not a weight of F).
\item
  Deduce that \(\mathrm{E}(\lambda)\otimes\mathrm{E}(\mu)\) contains a unique submodule isomorphic to
\end{enumerate}

\[
\mathrm{E} (\lambda + \mu - \alpha).
\]

\begin{enumerate}
\def\labelenumi{\alph{enumi})}
\setcounter{enumi}{3}
\tightlist
\item
  Show that \(\mathbf{S}^{2}(\mathrm{E}(\lambda))\) (resp. \(\bigwedge^{2}\mathrm{E}(\lambda)\) ) contains a unique submodule isomorphic to \(\mathrm{E}(2\lambda)\) (resp. to \(\mathrm{E}(2\lambda-\alpha)\) ).
\end{enumerate}

¶ 18) Choose a positive non-degenerate symmetric bilinear form \((\cdot|\cdot)\) on \(\mathfrak{h}_{\mathbf{R}}^{*}\) invariant under W. Let \(\lambda \in P_{++}\) .

\begin{enumerate}
\def\labelenumi{\alph{enumi})}
\item
  Let \(\mu\) be a weight of \(\mathrm{E}(\lambda)\) . Write \(\lambda - \mu\) in the form \(\sum_{\alpha \in \mathrm{B}} k_{\alpha} \alpha\) . Let \(\alpha \in \mathrm{B}\) be such that \(k_{\alpha} \neq 0\) . Show that there exists \(\alpha_{1}, \ldots, \alpha_{n} \in \mathrm{B}\) such that \((\lambda | \alpha_{1}) \neq 0\) , \((\alpha_{1} | \alpha_{2}) \neq 0, \ldots, (\alpha_{n-1} | \alpha_{n}) \neq 0\) , \((\alpha_{n} | \alpha) \neq 0\) .
\item
  Let \(v\) be a primitive element of \(\mathrm{E}(\lambda)\) , and let \(\alpha_1, \ldots, \alpha_n \in \mathrm{B}\) satisfy the following conditions:
\end{enumerate}

\begin{enumerate}
\def\labelenumi{(\roman{enumi})}
\item
  \((\alpha_{i}|\alpha_{i+1})\neq0\) for \(i=1,2,\ldots,n-1;\)
\item
  \((\alpha_{i}|\alpha_{j}) = 0\) for \(j > i + 1\) ;
\item
  \(\lambda(H_{\alpha_{1}})\neq0\) and \(\lambda(H_{\alpha_{2}})=\cdots=\lambda(H_{\alpha_{n}})=0.\)
\end{enumerate}

Show that \(X_{-\alpha_n}X_{-\alpha_{n-1}} \ldots X_{-\alpha_1}v \neq 0\) . (Observe that, for \(1 \leq s \leq n\) ,

\[
\lambda - \alpha_ {1} - \dots - \alpha_ {s - 1} + \alpha_ {n}
\]

is not a weight of \(\mathrm{E}(\lambda)\) , and deduce that \(X_{\alpha_s}X_{-\alpha_s}X_{-\alpha_{s-1}}\ldots X_{-\alpha_1}v\neq 0\) by induction on \(s\) .) If \(\sigma \in \mathfrak{S}_n\) and \(\sigma \neq 1\) , then \(X_{-\alpha_{\sigma(n)}}X_{-\alpha_{\sigma(n-1)}}\ldots X_{-\alpha_{\sigma(1)}}v = 0\) . (Let \(r\) be the smallest integer such that \(\sigma(r)\neq r\) . Use \(a\) ) to show that \(\lambda -\alpha_{\sigma(1)} - \ldots -\alpha_{\sigma(r)}\) is not a weight of \(\mathrm{E}(\lambda)\) .)

\begin{enumerate}
\def\labelenumi{\alph{enumi})}
\setcounter{enumi}{2}
\tightlist
\item
  Let \(\lambda' \in \mathrm{P}_{++}\) . A chain joining \(\lambda\) to \(\lambda'\) is a sequence \((\alpha_1, \ldots, \alpha_n)\) of elements of B such that \(n \geq 1\) , \((\lambda|\alpha_1) \neq 0\) , \((\alpha_1|\alpha_2) \neq 0, \ldots, (\alpha_{n-1}|\alpha_n) \neq 0\) , \((\alpha_n|\lambda') \neq 0\) . Such a chain is called minimal if no strict subsequence of \((\alpha_1, \ldots, \alpha_n)\) joins \(\lambda\) to \(\lambda'\) . In that case, we have \((\alpha_i|\alpha_j) = 0\) if \(|i - j| \geq 2\) , \((\lambda|\alpha_i) = 0\) if \(i \geq 2\) and \((\lambda'|\alpha_i) = 0\) if \(i \leq n - 1\) .
\end{enumerate}

\(d\) ) Let \((\alpha_{1},\ldots ,\alpha_{n})\) be a minimal chain joining \(\lambda\) to \(\lambda^{\prime}\) . If \(v^{\prime}\) is a primitive vector in \(\operatorname {E}(\lambda ')\) , put:

\[
v _ {s} = X _ {- \alpha_ {s}} X _ {- \alpha_ {s - 1}} \ldots X _ {- \alpha_ {1}} v (s = 0, 1, \ldots , n)
\]

\[
v _ {s} ^ {\prime} = X _ {- \alpha_ {s + 1}} X _ {- \alpha_ {s + 2}} \dots X _ {- \alpha_ {n}} v ^ {\prime} (s = 0, 1, \ldots , n)
\]

\(a_0 = (\lambda |\alpha_1), a_n = (-1)^n (\lambda '|\alpha_n), a_s = (-1)^{s + 1}(\alpha_s|\alpha_{s + 1}), 1 \leq s \leq n - 1.\) Show that

\[
\sum_ {s = 0} ^ {n} a _ {s} v _ {s} \otimes v _ {s} ^ {\prime}
\]

is a primitive element of \(\mathrm{E}(\lambda)\otimes\mathrm{E}(\lambda')\) of weight \(\lambda+\lambda'-\alpha_{1}-\cdots-\alpha_{n}\) , and that it is unique up to homothety.

(Use b) and c) to show that every element of \(\mathrm{E}(\lambda) \otimes \mathrm{E}(\lambda')\) of weight

\[
\lambda + \lambda^ {\prime} - \alpha_ {1} - \dots - \alpha_ {n}
\]

is a linear combination of the \(v_{s} \otimes v_{s}^{\prime}\) . Then write down the condition for such a linear combination to be a primitive vector.)

Deduce that \(\mathrm{E}(\lambda) \otimes \mathrm{E}(\lambda')\) contains a unique g-submodule isomorphic to

\[
\operatorname{E} \left(\lambda + \lambda^ {\prime} - \alpha_ {1} - \dots - \alpha_ {n}\right).
\]

(When \(n = 1\) we recover Exerc. 17.)

\begin{enumerate}
\def\labelenumi{\alph{enumi})}
\setcounter{enumi}{4}
\tightlist
\item
  Let w be a primitive element of \(\mathrm{E}(\lambda)\otimes\mathrm{E}(\lambda')\) . Assume that the weight \(\nu\) of w is distinct from \(\lambda+\lambda'\) . Show that there exists a chain \((\alpha_{1},\ldots,\alpha_{n})\) joining \(\lambda\) to \(\lambda'\) such that
\end{enumerate}

\[
\nu \leq \lambda + \lambda^ {\prime} - \alpha_ {1} - \dots - \alpha_ {n}.
\]

(Let C be the set of \(\alpha \in \mathrm{B}\) such that the coordinate of index \(\alpha\) of \(\lambda + \lambda' - \nu\) is \(\neq 0\) . Let D (resp. D') be the set of \(\alpha \in \mathrm{C}\) such that there exists \(\gamma_1, \ldots, \gamma_t \in \mathrm{C}\) satisfying \((\lambda|\gamma_1) \neq 0\) (resp. \((\lambda'|\gamma_1) \neq 0\) ), \((\gamma_1|\gamma_2) \neq 0, \ldots, (\gamma_{t-1}|\gamma_t) \neq 0\) , \((\gamma_t|\alpha) \neq 0\) . Let Y (resp. Y') be the set of weights of E( \(\lambda\) ) (resp. E( \(\lambda'\) )) of the form \(\lambda - \sum_{\alpha \in \mathrm{C}} k_\alpha\alpha\) (resp. \(\lambda' - \sum_{\alpha \in \mathrm{C}} k_\alpha\alpha\) ), with \(k_\alpha \in \mathbf{N}\) . Show that \(w\) belongs to \(\left(\sum_{\mu \in \mathrm{Y}}\mathrm{E}(\lambda)^{\mu}\right)\otimes \left(\sum_{\mu^{\prime}\in \mathrm{Y}^{\prime}}\mathrm{E}(\lambda^{\prime})^{\mu^{\prime}}\right)\) . Using \(a\) ) and the fact that \(\nu \neq \lambda +\lambda '\) , show that \(\mathrm{D}\cap \mathrm{D}'\neq \emptyset .\)

\(f\) ) Show the equivalence of the following properties:

\begin{enumerate}
\def\labelenumi{(\roman{enumi})}
\item
  \(\mathrm{E}(\lambda)\otimes \mathrm{E}(\lambda^{\prime})\) is isomorphic to \(\mathrm{E}(\lambda +\lambda^{\prime})\)
\item
  \(\mathrm{E}(\lambda) \otimes \mathrm{E}(\lambda')\) is a simple module.
\item
  There is no chain joining \(\lambda\) to \(\lambda'\) .
\item
  R is the direct sum of two root systems \(R_{1}\) and \(R_{1}^{\prime}\) such that \(\lambda \in \mathrm{P}(\mathrm{R}_{1})\) and \(\lambda^{\prime} \in \mathrm{P}(\mathrm{R}_{1}^{\prime})\) .
\item
  \(\mathfrak{g}\) is the product of two ideals \(\mathfrak{s}\) and \(\mathfrak{s}'\) such that \(\mathfrak{s}'\) .E(λ) = 0 and \(\mathfrak{s}\) .E(λ') = 0.
\end{enumerate}

(Use \(d\) ) to prove the equivalence of (ii) and (iii).) \(^{15}\)

\begin{enumerate}
\def\labelenumi{\arabic{enumi})}
\setcounter{enumi}{18}
\item
  We recall the notations of Prop. 10. Let \(\bar{k}\) be an algebraic closure of \(k\) . If \(x \in \mathfrak{g}\) , denote by \(\mathcal{X}_{\mathrm{E}}(x)\) (resp. \(\mathcal{X}_{\mathrm{F}}(x)\) , resp. \(\mathcal{X}_{\mathrm{G}}(x)\) ) the set of eigenvalues of \(x_{\mathrm{E}}\) (resp. \(x_{\mathrm{F}}\) , resp. \(x_{\mathrm{G}}\) ) in \(\bar{k}\) . Show that \(\mathcal{X}_{\mathrm{G}}(x) = \mathcal{X}_{\mathrm{E}}(x) + \mathcal{X}_{\mathrm{F}}(x)\) for all \(x \in \mathfrak{g}\) and that, when E and F are given, this property characterizes the simple \(\mathfrak{g}\) -module G up to isomorphism.
\item
  Assume that \(k = \mathbf{R}\) or \(\mathbf{C}\) . Let \(\Gamma\) be a simply-connected Lie group with Lie algebra \(\mathfrak{g}\) . Let \(\lambda, \mu, \mathrm{E}, \mathrm{F}, \mathrm{G}\) be as in Prop. 10. Let \((e_1, \ldots, e_n)\) (resp. \((f_1, \ldots, f_p)\) ) be a basis of \(\mathrm{E}\) (resp. \(\mathrm{F}\) ) consisting of eigenvectors of \(\mathfrak{h}\) , with \(e_1 \in \mathrm{E}^\lambda\) and \(f_1 \in \mathrm{F}^\mu\) . We can consider \(\mathrm{E}, \mathrm{F}, \mathrm{G}\) as \(\Gamma\) -modules. If \(\gamma \in \Gamma\) , denote by \(a_i(\gamma)\) the coordinate of \(\gamma.e_1\) with index \(i\) , and \(b_j(\gamma)\) the coordinate of \(\gamma.f_1\) with index \(j\) . Show that the function \(a_i b_j\) on \(\Gamma\) is not identically zero. Deduce that, for all \((i,j)\) , there exists an element of \(\mathrm{G} \subset \mathrm{E} \otimes \mathrm{F}\) whose coordinate of index \((i,j)\) is \(\neq 0\) ; for \(k = \mathbf{R}\) or \(\mathbf{C}\) , this gives a new proof of Prop. 10. Pass from this to the case \(k = \mathbf{Q}\) , and then to the case of an arbitrary field, cf.~Exerc. 4.
\item
  Let \(\lambda, \mu \in \mathrm{P}_{++}\) . Let E,F,G be simple \(\mathfrak{g}\) -modules of highest weights \(\lambda, \mu, \lambda + \mu\) , and let \(n\) be an integer \(\geq 1\) . If \(\omega\) is a weight of E of multiplicity \(n\) , then \(\omega + \mu\) is a weight of G of multiplicity \(\geq n\) . (We have \(\mathrm{G}^{\omega + \mu} \subset \bigoplus_{\nu + \sigma = \omega + \mu} \mathrm{E}^{\nu} \otimes \mathrm{F}^{\sigma}\) . If \(\dim \mathrm{G}^{\omega + \mu} < n\) , the projection of \(\mathrm{G}^{\omega + \mu}\) onto \(\mathrm{E}^{\omega} \otimes \mathrm{F}^{\mu}\) is of the form \(\mathrm{E}' \otimes \mathrm{F}^{\mu}\) , with \(\mathrm{E}'\) strictly contained in \(\mathrm{E}^{\omega}\) . Derive a contradiction from this by choosing adapted bases of E and F and imitating the proof of Prop. 10.)
\item
  Assume that \(\mathfrak{g}\) is simple, in other words that R is irreducible. Show that there exists a unique dominant weight \(\lambda \neq 0\) such that the set of weights of \(\mathrm{E}(\lambda)\) is \(\mathrm{W}.\lambda \cup \{0\}\) : we have \(\lambda = \alpha\) , where \(\alpha \in \mathrm{R}\) is such that \(H_{\alpha}\) is the highest root of \(\mathrm{R}^{\vee}\) . When all the roots are of the same length (cases \(\mathrm{A}_l,\mathrm{D}_l,\mathrm{E}_6,\mathrm{E}_7,\mathrm{E}_8\) ), we have \(\lambda = \tilde{\alpha}\) ; this is the only root that is a dominant weight; the corresponding representation is the adjoint representation of g. In the other cases, \(\lambda\) is the only root of minimum length that is a dominant weight; with the notations of Chap. VI, Plates, we have: \(\lambda = \varpi_{1}\) (type \(B_{l}\) ); \(\lambda = \varpi_{2}\) (type \(C_{l}\) ); \(\lambda = \varpi_{4}\) (type \(F_{4}\) ); \(\lambda = \varpi_{1}\) (type \(G_{2}\) ).
\item
  Let \(\mathrm{U}^0\) be the commutant of \(\mathfrak{h}\) in \(\mathrm{U}(\mathfrak{g})\) (cf.~§6, no. 4).
\end{enumerate}

\begin{enumerate}
\def\labelenumi{\alph{enumi})}
\item
  Let V be a (finite dimensional) g-module. Show that the \(V^{\lambda}\) , \(\lambda \in P\) , are stable under \(U^{0}\) , and that, if V is simple and \(V^{\lambda} \neq 0\) , \(V^{\lambda}\) is a simple \(U^{0}\) -module (use the decomposition \(\mathrm{U}(\mathfrak{g}) = \oplus\mathrm{U}^{\lambda}\) , loc. cit.).
\item
  Show that, for every element \(c \neq 0\) of \(U^{0}\) , there exists a finite dimensional simple representation \(\rho\) of \(U^{0}\) such that \(\rho(c) \neq 0\) . (Use a) and Chap. I, §7, Exerc. 3 a).
\end{enumerate}

\begin{enumerate}
\def\labelenumi{\arabic{enumi})}
\setcounter{enumi}{23}
\tightlist
\item
  Let \(A\) be a unital associative algebra and \(M\) the set of its finite codimensional two-sided ideals.
\end{enumerate}

\begin{enumerate}
\def\labelenumi{\alph{enumi})}
\tightlist
\item
  Let \(\mathrm{U}^*\) be the vector space dual of U. If \(\theta \in \mathrm{U}^*\) , show the equivalence of the following properties:
\end{enumerate}

\begin{enumerate}
\def\labelenumi{(\roman{enumi})}
\item
  there exists \(\mathfrak{m} \in \mathrm{M}\) such that \(\theta(\mathfrak{m}) = 0\) ;
\item
  there exist two finite families \((\theta_{i}^{\prime})\) and \((\theta_{i}^{\prime\prime})\) of elements of \(U^{*}\) such that
\end{enumerate}

\[
\theta (x y) = \sum_ {i} \theta_ {i} ^ {\prime} (x) \theta_ {i} ^ {\prime \prime} (y) \quad \text { for   all } x, y \in \mathrm{U}.
\]

The elements \(\theta\) with these properties form a subspace \(\mathrm{U}'\) of \(\mathrm{U}^*\) , which coincides with that denoted by \(\mathrm{B}'\) in Algebra, Chap. III, §11, Exerc. 27. There exists a unique coalgebra structure on \(\mathrm{U}'\) whose coproduct \(c: \mathrm{U}' \to \mathrm{U}' \otimes \mathrm{U}'\) is given by

\[
c (\theta) = \sum_ {i} \theta_ {i} ^ {\prime} \otimes \theta_ {i} ^ {\prime \prime},
\]

where \(\theta_i', \theta_i'' \in \mathrm{U}'\) are such that \(\theta(xy) = \sum_{i} \theta_i'(x) \theta_i''(y)\) for all \(x, y \in \mathrm{U}\) , cf.~(ii).

The coalgebra \(\mathrm{U}'\) is the union of the increasing filtration by the subspaces \((\mathrm{U} / \mathfrak{m})^*\) , \(\mathfrak{m} \in \mathrm{M}\) , which are finite dimensional. The dual of \(\mathrm{U}'\) can be identified with the algebra \(\hat{\mathrm{U}} = \varprojlim \mathrm{U} / \mathfrak{m}\) ; if \(\hat{\mathrm{U}}\) is given the projective limit of the discrete topologies on the \(\mathrm{U} / \mathfrak{m}, \mathfrak{m} \in \mathrm{M}\) , the continuous linear forms on \(\hat{\mathrm{U}}\) are given by the elements of \(\mathrm{U}'\) .


\begin{enumerate}
\def\labelenumi{\alph{enumi})}
\setcounter{enumi}{1}
\tightlist
\item
  Let E be a finite dimensional left U-module. Its annihilator \(m_{E}\) belongs to M; the composite \(\hat{U} \to U/m_{E} \to \text{End}(E)\) gives E the structure of a left \(\hat{U}\) -module. If F is a finite dimensional left U-module, a linear map \(f : E \to F\) is a U-homomorphism if and only if it is a \(\hat{U}\) -homomorphism. If \(a \in E, b \in E^{*}\) , the linear form \(\theta_{a,b} : x \mapsto \langle xa, b \rangle\) belongs to \(U'\) , and
\end{enumerate}

\[
\langle x, \theta_ {a, b} \rangle = \langle x a, b \rangle \quad \text { for   all } x \in \hat {\mathrm{U}}.
\]

The \(\theta_{a,b}\) (for varying E, \(a,b\) ) generate the \(k\) -vector space U'.

\begin{enumerate}
\def\labelenumi{\alph{enumi})}
\setcounter{enumi}{2}
\tightlist
\item
  Let \(X_{U}\) be the set of isomorphism classes of finite dimensional left U-modules. For all \(E \in X_{U}\) , let \(u_{E}\) be a k-linear endomorphism of E; assume that
\end{enumerate}

\[
f \circ u _ {\mathrm{E}} = u _ {\mathrm{F}} \circ f \quad \text { for   all   E, F } \in X_{U} \text { and } f \in \operatorname{Hom} _ {\mathrm{U}} (\mathrm{E,F}).
\]

Show that there exists a unique element \(x \in \hat{U}\) such that \(x_{E} = u_{E}\) for all \(E \in X_{U}\) . (Reduce to the case in which U is finite dimensional.)

¶ 25) Let a be a Lie algebra and U its enveloping algebra. Apply the definitions and results of Exerc. 24 to the algebra U. In particular, \(U' \subset U^*\) and the dual of \(U'\) can be identified with the algebra \(\hat{U} = \varprojlim U/m\) ; the canonical map \(U \to \hat{U}\) is injective (Chap. I, §7, Exerc. 3).

\begin{enumerate}
\def\labelenumi{\alph{enumi})}
\tightlist
\item
  The coalgebra structure of U (Chap. II, §1, no. 4) defines an algebra structure on the dual U* (cf.~Chap. II, §1, no. 5, Prop. 10, and Algebra, Chap. III, §11, no. 2). If E, F are finite dimensional a-modules,
\end{enumerate}

\[
\theta_ {a, b} \circ \theta_ {c, d} = \theta_ {a \otimes c, b \otimes d} \text { for } a \in \mathrm{E}, b \in \mathrm{E} ^ {*}, c \in \mathrm{F}, d \in \mathrm{F} ^ {*}.
\]

Deduce that \(U'\) is a subalgebra of \(U^{*}\) . The coalgebra and algebra structures of \(U'\) make it a commutative bigebra (Algebra, Chap. III, §11, no. 4).

\begin{enumerate}
\def\labelenumi{\alph{enumi})}
\setcounter{enumi}{1}
\tightlist
\item
  Let \(x\) be an element of \(\hat{\mathrm{U}}\) ; identify \(x\) with a linear form \(\mathrm{U}' \to k\) . Prove the equivalence of the following properties:
\end{enumerate}

\begin{enumerate}
\def\labelenumi{(\roman{enumi})}
\item
  \(x\) is a homomorphism of algebras from \(\mathrm{U}'\) to \(k\) .
\item
  \(x_{E \otimes F} = x_{E} \otimes x_{F}\) for all finite dimensional a-modules E and F.
\end{enumerate}

(Prove first that (ii) is equivalent to

\[
\left(\mathrm{ii} ^ {\prime}\right) x \left(\theta_ {a \otimes c, b \otimes d}\right) = x \left(\theta_ {a, b}\right) x \left(\theta_ {c, d}\right) \text {if} a \in \mathrm{E}, b \in \mathrm{E} ^ {*}, c \in \mathrm{F}, d \in \mathrm{F} ^ {*},
\]

and use the fact that the \(\theta_{a,b}\) generate \(\mathrm{U}'\) .)

\(c)\) Let \(x\) be an element of \(\hat{\mathbf{U}}\) satisfying conditions (i) and (ii) of \(b\) ). Show the equivalence of the following conditions:

\begin{enumerate}
\def\labelenumi{(\roman{enumi})}
\setcounter{enumi}{2}
\item
  \(x\) takes the unit element of \(\mathrm{U}'\) to the unit element of \(k\) .
\item
  \(x \neq 0.\)
\item
  If \(k\) is given the trivial \(\mathfrak{a}\) -module structure, we have \(x_{k} = \mathrm{Id}\) .
\item
  \(x_{\mathrm{E}}\) is invertible for all E.
\end{enumerate}

(The equivalence (iii) \(\Longleftrightarrow\) (iv) follows from the fact that \(x\) is a homomorphism of algebras. On the other hand, \(x_{k} = \lambda \mathrm{Id}\) , with \(\lambda \in k\) . Using the \(\mathfrak{a}\) -isomorphism \(k \otimes \mathrm{E} \to \mathrm{E}\) , deduce that \(\lambda x_{\mathrm{E}} = x_{\mathrm{E}}\) for all E, and, in particular, that \(\lambda^2 = \lambda\) by taking \(\mathrm{E} = k\) . The case \(\lambda = 1\) corresponds to \(x \neq 0\) , hence (iv) \(\Longleftrightarrow\) (v), and (vi) \(\Longrightarrow\) (v). To prove that (v) \(\Longrightarrow\) (vi), show that, if F is the dual of E, then \(^t x_{\mathrm{F}} \circ x_{\mathrm{E}} = \lambda \mathrm{Id}_{\mathrm{E}}\) .)

\begin{enumerate}
\def\labelenumi{\alph{enumi})}
\setcounter{enumi}{3}
\tightlist
\item
  Let G be the set of elements of \(\hat{U}\) satisfying conditions (i) to (vi) above. Show that G is a subgroup of the group of invertible elements of \(\hat{U}\) .
\end{enumerate}

Let \(x \in G\) . If E is a finite dimensional \(\mathfrak{a}\) -module, then \(x_{\mathrm{E}} \in \mathbf{GL}(\mathrm{E})\) . This applies in particular to \(\mathrm{E} = \mathfrak{a}\) , equipped with the adjoint representation; this gives an element \(x_{\mathfrak{a}} \in \mathbf{GL}(\mathfrak{a})\) . Show that \(x_{\mathfrak{a}}\) is an automorphism of \(\mathfrak{a}\) (use the \(\mathfrak{a}\) -homomorphism \(\mathfrak{a} \otimes \mathfrak{a} \to \mathfrak{a}\) given by the bracket). This gives a homomorphism \(v: \mathrm{G} \to \operatorname{Aut}(\mathfrak{a})\) . If \(\mathrm{E}\) is a finite dimensional \(\mathfrak{a}\) -module,

\[
x _ {\mathrm{E}} (y. e) = v (x) (y). x _ {\mathrm{E}} (e) \quad \mathrm{if} y \in \mathfrak {a}, e \in \mathrm{E}
\]

(use the \(\mathfrak{a}\) -homomorphism \(\mathfrak{a} \otimes \mathrm{E} \to \mathrm{E}\) given by the operation of \(\mathfrak{a}\) on \(\mathrm{E}\) ).

\begin{enumerate}
\def\labelenumi{\alph{enumi})}
\setcounter{enumi}{4}
\tightlist
\item
  The principal anti-automorphism \(\sigma\) of U extends by continuity to \(\hat{\mathbf{U}}\) . Its transpose leaves \(\mathrm{U}'\) stable and induces an inversion on \(\mathrm{U}'\) (Algebra, Chap. III, §11, Exerc. 4). If \(x \in \mathrm{G}\) , then \(\sigma(x) = x^{-1}\) .
\end{enumerate}

\(f\) ) Assume that \(\mathfrak{a}\) is semi-simple \(^{16}\) . Let \(n\) be a nilpotent element of \(\mathfrak{a}\) . Then there exists a unique element \(e^n\) of \(\mathrm{G}\) such that \((e^n)_{\mathrm{E}} = \exp (n_{\mathrm{E}})\) for every finite dimensional \(\mathfrak{a}\) -module E. We have \(v(e^n) = \exp (\operatorname {ad}n)\in \operatorname {Aut}(\mathfrak{a})\) , so \(\operatorname {Aut}_e(\mathfrak{a})\subset v(\mathrm{G})\) .

Show that, if b is a subalgebra of a consisting of nilpotent elements, then

\[
e ^ {n}. e ^ {m} = e ^ {\mathrm{H} (n, m)} \quad \text { for } n, m \in \mathfrak {b},
\]

where H denotes the Hausdorff series (Chap. II, §6).

¶ 26) Apply the notations and results of Exerc. 25 to the case in which a = g (split case).

\begin{enumerate}
\def\labelenumi{\alph{enumi})}
\item
  Let \(x \in G\) and let \(\sigma = v(x)\) be its image in \(\operatorname{Aut}(\mathfrak{g})\) . If \(\rho\) is a representation of g, \(\rho\) and \(\rho \circ \sigma\) are equivalent. Deduce (cf.~no. 2, Remark 1) that \(\sigma\) belongs to \(\operatorname{Aut}_{0}(\mathfrak{g})\) . Extend this result to arbitrary semi-simple algebras.
\item
  Let \(\varphi \in \mathrm{T}_{\mathrm{P}} = \mathrm{Hom}(\mathrm{P}, k^{*})\) , where \(\mathrm{P} = \mathrm{P}(\mathrm{R})\) . If \(\mathrm{E}\) is a \(\mathfrak{g}\) -module, let \(\varphi_{\mathrm{E}}\) be the endomorphism of \(\mathrm{E}\) whose restriction to each \(\mathrm{E}^{\lambda}\) ( \(\lambda \in \mathrm{P}\) ) is the homothety of ratio \(\varphi(\lambda)\) . Show that there exists a unique element \(t(\varphi) \in \mathrm{G}\) such that \(t(\varphi)_{\mathrm{E}} = \varphi_{\mathrm{E}}\) for all \(\mathrm{E}\) (Use Exerc. 24 c), and the characterizations (ii) and (vi) of Exerc. 25.) This gives a homomorphism \(t: \mathrm{T}_{\mathrm{P}} \to \mathrm{G}\) . Show that \(t\) is injective. We use this to identify \(\mathrm{T}_{\mathrm{P}}\) with a subgroup of \(\mathrm{G}\) . The composite \(\mathrm{T}_{\mathrm{P}} \to \mathrm{G} \to \mathrm{Aut}(\mathfrak{g})\) is the homomorphism denoted by \(f \circ q\) in §5, no. 2.
\item
  Let \(x \in G\) be such that \(\sigma = v(x)\) belongs to the subgroup \(f(\mathrm{T}_{\mathrm{Q}})\) of \(\mathrm{Aut}_{0}(\mathfrak{g})\) (§5, no. 2), in other words, such that it operates trivially on h; denote the element of \(T_{Q} = \mathrm{Hom}(Q, k^{*})\) corresponding to x by \(\psi\) . We are going to show that x belongs to \(T_{P}\) . Prove successively:
\end{enumerate}

\(c_{1}\) ) If E is a \(\mathfrak{g}\) -module, \(x_{\mathrm{E}}\) is an \(\mathfrak{h}\) -endomorphism of E.

(Use the g-homomorphism \(g \otimes E \to E\) , and the fact that x operates trivially on h.) In particular, the \(E^{\mu}\) are stable under \(x_{E}\) .

\(c_{2}\) ) There exists \(\varphi \in \mathrm{T}_{\mathrm{P}}\) such that, for every \(\mathfrak{g}\) -module E, and every primitive element \(e\) of E of weight \(\lambda\) , we have \(x_{\mathrm{E}}e = \varphi (\lambda)e\) .

(Choose \(\varphi\) such that this relation is true when E is a fundamental module \(\mathrm{E}(\varpi_{\alpha})\) . Deduce the case of the \(\mathrm{E}(\lambda)\) , \(\lambda \in \mathrm{P}_{++}\) , by using the embedding of such a module in a tensor product of the \(\mathrm{E}(\varpi_{\alpha})\) . Pass from this to the general case.)

\(c_{3}\) ) Choose \(\varphi\) as in \(c_{2}\) ). Let E be a simple g-module of highest weight \(\lambda\) , and let \(\mu\) be a weight of E; then \(\lambda - \mu \in Q\) . Show that the restriction of \(x_{E}\) to \(E^{\mu}\) is the homothety of ratio \(\varphi(\lambda)\psi(\mu - \lambda)\) . (Same method as for \(c_{1}\) .)

\(c_{4}\) ) If \(\lambda, \mu \in P_{++}\) , and if \(\alpha \in B\) is not orthogonal to either \(\lambda\) or \(\mu\) , then

\[
\varphi (\lambda + \mu - \alpha) = \varphi (\lambda + \mu) \psi (- \alpha),
\]

so \(\varphi (\alpha) = \psi (\alpha)\) . (Use \(c_{2}\) ), \(c_{3}\) ), and the embedding of \(\operatorname {E}(\lambda +\mu -\alpha)\) in \(\operatorname {E}(\lambda)\otimes \operatorname {E}(\mu)\) , cf.~Exerc. 17.)

\(c_{5})\) Deduce from \(c_{4})\) that \(\varphi|\mathbf{Q}=\psi\) , and use \(c_{3})\) to deduce that \(x=t(\varphi)\) . \(d)\) Identify \(\mathrm{T}_{\mathrm{Q}}\) with a subgroup of \(\mathrm{Aut}_{0}(\mathfrak{g})\) by means of \(f\) . By \(a)\) , we have

\[
\operatorname{Aut} _ {e} (\mathfrak {g}) \subset v (\mathrm{G}) \subset \operatorname{Aut} _ {0} (\mathfrak {g})
\]

and, by \(c\) ), \(v(\mathrm{G}) \cap \mathrm{T}_{\mathrm{Q}} = \mathrm{Im}(\mathrm{T}_{\mathrm{P}})\) . Deduce (cf.~§5, no. 3) that \(\mathrm{Aut}_e(\mathfrak{g}) \cap \mathrm{T}_{\mathrm{Q}} = \mathrm{Im}(\mathrm{T}_{\mathrm{P}})\) and that \(f(\mathrm{G}) = \mathrm{Aut}_e(\mathfrak{g})\) . The canonical map

\[
\iota : \mathrm{T} _ {\mathrm{Q}} / \operatorname{Im} (\mathrm{T} _ {\mathrm{P}}) \to \operatorname{Aut} _ {0} (\mathfrak {g}) / \operatorname{Aut} _ {e} (\mathfrak {g})
\]

is therefore an isomorphism.

\begin{enumerate}
\def\labelenumi{\alph{enumi})}
\setcounter{enumi}{4}
\tightlist
\item
  The kernel of \(v: \mathrm{G} \to \mathrm{Aut}_e(\mathfrak{g})\) is equal to the kernel of \(\mathrm{T}_{\mathrm{P}} \to \mathrm{T}_{\mathrm{Q}}\) ; it is isomorphic to
\end{enumerate}

\[
\operatorname{Hom} (\mathrm{P} / \mathrm{Q}, k ^ {*});
\]

this is a finite abelian group contained in the centre of G, and its order divides \((P:Q)\) ; if k is algebraically closed, it is isomorphic to the dual of P/Q (Algebra, Chap. VII, §4, no. 8).

\(f)\) Let \(\alpha \in \mathrm{R}\) , \(X_{\alpha} \in \mathfrak{g}^{\alpha}\) and \(X_{-\alpha} \in \mathfrak{g}^{-\alpha}\) be such that \([X_{\alpha}, X_{-\alpha}] = -H_{\alpha}\) , and let \(\rho_{\alpha}\) be the corresponding representation of \(\mathfrak{sl}(2, k)\) on \(\mathfrak{g}\) . If E is a \(\mathfrak{g}\) -module, deduce (§1, no. 4) a representation of \(\mathbf{SL}(2, k)\) on E, and hence (Exerc. 25 b), c)) a homomorphism

\[
\varphi_ {\alpha}: \mathbf {SL} (2, k) \to \mathrm{G}.
\]

Show that \(\operatorname{Im}(\varphi_{\alpha})\) contains the elements of \(T_{P}\) of the form \(\lambda \mapsto t^{\lambda(H_{\alpha})}, t \in k^{*}\) . Deduce that the \(\operatorname{Im}(\varphi_{\alpha}), \alpha \in B\) , generate G (show first that the group they generate contains \(T_{P}\) ). In particular, G is generated by the \(e^{n}\) , with \(n \in g^{\alpha}, \alpha \in B \cup -B\) . The derived group of G is equal to G.

\(g\) ) If a subgroup \(\mathbf{G}'\) of \(\mathbf{G}\) is such that \(v(\mathrm{G}') = \mathrm{Aut}_e(\mathfrak{g})\) , then \(\mathrm{G}' = \mathrm{G}\) (use \(f\) ). \(h\) ) Let \(\mathrm{E}\) be a faithful \(\mathfrak{g}\) -module, and let \(\Gamma\) be the subgroup of \(\mathrm{P}\) generated by the weights of \(\mathrm{E}\) . Then \(\mathrm{P} \supset \Gamma \supset \mathrm{Q}\) , cf.~Exerc. 5. Show that the kernel of the canonical homomorphism \(\mathrm{G} \to \mathbf{GL}(\mathrm{E})\) is equal to the subgroup of \(T_{P}\) consisting of the elements \(\varphi\) whose restriction to \(\varGamma\) is trivial. In particular, if \(\varGamma = P\) the homomorphism \(\mathrm{G} \to \mathbf{GL}(\mathrm{E})\) is injective. If \(\varGamma = Q\) , this homomorphism factorizes as \(\mathrm{G} \stackrel{v}{\longrightarrow} \operatorname{Aut}_{e}(\mathfrak{g}) \longrightarrow \mathbf{GL}(\mathrm{E})\) , and the homomorphism \(\operatorname{Aut}_{e}(\mathfrak{g}) \to \mathbf{GL}(\mathrm{E})\) is injective.

\begin{enumerate}
\def\labelenumi{\arabic{enumi})}
\setcounter{enumi}{26}
\tightlist
\item
  Let \(\Omega = \mathrm{P} / \mathrm{Q}\) . If \(\omega \in \Omega\) , and if \(\mathrm{E}\) is a \(\mathfrak{g}\) -module, denote by \(\mathrm{E}_{\omega}\) the direct sum of the \(\mathrm{E}^{\lambda}\) , for \(\lambda \in \omega\) . We have \(\mathrm{E} = \bigoplus_{\omega \in \Omega} \mathrm{E}_{\omega}\) .
\end{enumerate}

\begin{enumerate}
\def\labelenumi{\alph{enumi})}
\tightlist
\item
  Show that \(E_{\omega}\) is a g-submodule of E. We have \((\mathrm{E}^{*})_{\omega} = (\mathrm{E}_{-\omega})^{*}\) and
\end{enumerate}

\[
(\mathrm{E} \otimes \mathrm{F}) _ {\omega} = \bigoplus_ {\alpha + \beta = \omega} \mathrm{E} _ {\alpha} \otimes \mathrm{F} _ {\beta}
\]

if F is another g-module.

\begin{enumerate}
\def\labelenumi{\alph{enumi})}
\setcounter{enumi}{1}
\item
  Let \(\chi\in\mathrm{Hom}(\varOmega,k^{*})=\mathrm{Ker}(\mathrm{T}_{\mathrm{P}}\to\mathrm{T}_{\mathrm{Q}})\) . Identify \(\chi\) with an element of the kernel of \(f:\mathrm{G}\to\mathrm{Aut}_{e}(\mathfrak{g})\) , cf.~Exerc. 26 e). Show that the operation of \(\chi\) on \(E_{\omega}\) is the homothety of ratio \(\chi(\omega)\) .
\item
  What are the \(\mathrm{E}_{\omega}\) when \(\mathfrak{g} = \mathfrak{sl}(2,k)\) ?
\end{enumerate}

\exercisesection{\S~8}

\begin{enumerate}
\def\labelenumi{\arabic{enumi})}
\item
  Let \(f\) be an invariant polynomial function on \(\mathfrak{g}\) . Show that \(f\) is invariant under \(\operatorname{Aut}(\mathfrak{g})\) if and only if \(f|\mathfrak{h}\) is invariant under \(\operatorname{Aut}(\mathrm{R})\) . Deduce that, if the Dynkin graph of \(\mathrm{R}\) has a non-trivial automorphism, there exists an invariant polynomial function on \(\mathfrak{g}\) that is not invariant under \(\operatorname{Aut}(\mathfrak{g})\) .
\item
  Take \(\mathfrak{g} = \mathfrak{sl}(3,k)\) . Show that \(x\mapsto \det (x)\) is an invariant polynomial function on \(\mathfrak{g}\) that is not invariant under \(\operatorname {Aut}(\mathfrak{g})\) (use the automorphism \(x\mapsto -^t x\) ).
\item
  Let \(\mathfrak{a}\) be a semi-simple Lie algebra, and \(s\in \mathrm{Aut}(\mathfrak{a})\) . Show the equivalence of:
\end{enumerate}

\begin{enumerate}
\def\labelenumi{(\roman{enumi})}
\item
  \(s \in \text{Aut}_{0}(\mathfrak{a})\) .
\item
  \(s\) operates trivially on the centre of \(\mathrm{U}(\mathfrak{a})\) .
\item
  For all \(x \in \mathfrak{a}\) , there exists \(t \in \mathrm{Aut}_0(\mathfrak{a})\) such that \(tx = sx\) .
\end{enumerate}

(Use Prop. 6 to show that (iii) \(\Longrightarrow\) (i).)

\begin{enumerate}
\def\labelenumi{\arabic{enumi})}
\setcounter{enumi}{3}
\item
  Show that, in Cor. 2 of Prop. 2, and in Th. 1 (ii), we can restrict ourselves to representations \(\rho\) whose weights are radical weights (remark that Prop. 1 remains valid when \(k[\mathrm{P}]^{\mathrm{W}}\) is replaced by \(k[\mathrm{Q}]^{\mathrm{W}}\) , where \(\mathrm{Q}\) is the group of radical weights).
\item
  We retain the notations of §6 and §7. Let \(\lambda \in \mathfrak{h}^*\) . If, for any \(w \in \mathrm{W}\) , \(w \neq 1\) , we have \((\lambda + \rho) - w(\lambda + \rho) \notin \mathrm{Q}_+\) , then \(\mathrm{Z}(\lambda)\) is simple.
\end{enumerate}

(Use Cor. 1 (ii) of Th. 2.)

\begin{enumerate}
\def\labelenumi{\arabic{enumi})}
\setcounter{enumi}{5}
\item
  Let \(\mathfrak{a}\) be a Lie algebra, \(f\) a polynomial function on \(\mathfrak{a}\) , and \(x, y\) two elements of \(\mathfrak{a}\) . Put \(f_y = \theta^*(y)f\) (cf.~no. 3), and denote by \(\mathrm{D}_x f\) the tangent linear map of \(f\) at \(x\) (Chap. VII, App. I, no. 2). Show that \(f_y(x) = (\mathrm{D}_xf)([x,y])\) . Deduce that \(\mathrm{D}_xf\) vanishes on \(\operatorname{Im}(\operatorname{ad}x)\) when \(f\) is invariant.
\item
  *Let \(d_1, \ldots, d_l\) be the characteristic degrees of the algebra I(g), cf.~Chap. V, §5, no. 1. For any integer \(n \geq 0\) , denote by \(r_n\) the number of elements of degree \(n\) in a homogeneous basis of S(g) over I(g) (cf.~no. 3, Remark 2), and put \(r(\mathrm{T}) = \sum_{n=0}^{\infty} r_n \mathrm{T}^n\) . Show that
\end{enumerate}

\[
r (\mathrm{T}) = (1 - \mathrm{T}) ^ {- \mathrm{N}} \prod_ {i = 1} ^ {l} (1 - \mathrm{T} ^ {d _ {i}}), \quad \text { where } \mathrm{N} = \dim (\mathfrak {g}).
\]

\begin{enumerate}
\def\labelenumi{\arabic{enumi})}
\setcounter{enumi}{7}
\tightlist
\item
  Put \(l = \operatorname{rk}(\mathfrak{g}), N = \dim(\mathfrak{g})\) . If \(x \in \mathfrak{g}\) , define \(a_i(x)\) , \(0 \leq i \leq N\) , by the formula
\end{enumerate}

\[
\det (\mathrm{T} + \operatorname{ad} x) = \sum_ {i = 0} ^ {\mathrm{N}} \mathrm{T} ^ {\mathrm{N} - i} a _ {i} (x).
\]

The function \(a_{i}\) thus defined is homogeneous polynomial of degree i, and invariant under \(\operatorname{Aut}(\mathfrak{g})\) . If \(x \in h\) and \(i \leq N - l\) , \(a_{i}(x)\) is the ith elementary symmetric function of the \(\alpha(x)\) , \(\alpha \in R\) ; in particular, \(a_{N-l}(x) = \prod_{\alpha \in \mathbb{R}} \alpha(x)\) . Construct an example in which the \(a_{i}\) do not generate the algebra of polynomial functions on g invariant under \(\operatorname{Aut}(\mathfrak{g})\) .

\begin{enumerate}
\def\labelenumi{\arabic{enumi})}
\setcounter{enumi}{8}
\tightlist
\item
  The notations are those of no. 5. Let \(\lambda \in \mathfrak{h}^*\) , \(z \in \mathbf{Z}\) , \(z'\) the image of \(z\) under the principal anti-automorphism of \(\mathrm{U}(\mathfrak{g})\) , and \(w_0\) the element of W that transforms B into \(-\mathrm{B}\) . Show that \(\chi_{\lambda}(z) = \chi_{-w_0\lambda}(z')\) .
\end{enumerate}

(It suffices to prove this for \(\lambda \in \mathrm{P}_{++}\) . Consider the operation of \(z\) on \(\mathrm{E}(\lambda)\) and \(\mathrm{E}(\lambda)^*\) ; use Prop. 11 of §7.)

\begin{enumerate}
\def\labelenumi{\arabic{enumi})}
\setcounter{enumi}{9}
\tightlist
\item
  (In this exercise, and in the following three, we retain the notations of §6.) Let \(\lambda \in \mathfrak{h}^*\) .
\end{enumerate}

\begin{enumerate}
\def\labelenumi{\alph{enumi})}
\tightlist
\item
  Let \(\mathrm{N}, \mathrm{N}'\) be \(\mathfrak{g}\) -submodules of \(\mathrm{Z}(\lambda)\) such that \(\mathrm{N}' \subset \mathrm{N}\) and \(\mathrm{N} / \mathrm{N}'\) is simple. Show that there exists \(\mu \in \lambda - \mathrm{Q}_+\) such that \(\mathrm{N} / \mathrm{N}'\) is isomorphic to \(\mathrm{E}(\mu)\) (apply Th. 1 of §6), and such that \(\mu + \rho \in \mathrm{W}.(\lambda + \rho)\) (apply Cor. 1 of Th. 2).
\item
  Show that \(\mathrm{Z}(\lambda)\) admits a Jordan-Hölder sequence. (Apply a) and the fact that the weights of \(\mathrm{Z}(\lambda)\) are of finite multiplicity.)
\end{enumerate}

\begin{enumerate}
\def\labelenumi{\arabic{enumi})}
\setcounter{enumi}{10}
\tightlist
\item
  Let \(\lambda \in \mathfrak{h}^*\) and let \(V\) be a non-zero \(\mathfrak{g}\) -submodule of \(Z(\lambda)\) . Show that there exists \(\mu \in \mathfrak{h}^*\) such that \(V\) contains a simple \(\mathfrak{g}\) -submodule isomorphic to \(Z(\mu)\) .
\end{enumerate}

(Let A be the set of \(\nu \in \mathfrak{h}^*\) such that V contains a \(\mathfrak{g}\) -submodule isomorphic to \(Z(\nu)\) . By using Prop. 6 of no. 6, show first that \(A \neq \varnothing\) . Then show that A is finite, and consider an element \(\mu\) of A such that \((\mu - Q_{+}) \cap A = \{\mu\}\) .)

¶ 12) a) Let \(\lambda, \mu \in \mathfrak{h}^*\) . Every non-zero \(\mathfrak{g}\) -homomorphism from \(Z(\mu)\) to \(Z(\lambda)\) is injective. (Use Prop. 6 of no. 6.)

\begin{enumerate}
\def\labelenumi{\alph{enumi})}
\setcounter{enumi}{1}
\item
  Let \(r \in \mathbf{N}\) , A a finite subset of \(\mathbf{N}^r\) , \(m = \mathrm{Card}(\mathrm{A})\) . For all \(\xi \in \mathbf{N}^r\) , let \(s(\xi)\) be the sum of the coordinates of \(\xi\) , and \(\mathfrak{P}_{\mathrm{A}}(\xi)\) the number of families \((n_{\alpha})_{\alpha \in \mathrm{A}}\) of integers \(\geq 0\) such that \(\xi = \sum_{\alpha \in \mathrm{A}} n_{\alpha}\alpha\) . Then \(\mathfrak{P}_{\mathrm{A}}(\xi) \leq (s(\xi) + 1)^m\) for all \(\xi \in \mathbf{N}^r\) . (Argue by induction on \(m\) .)
\item
  Let \(\lambda, \mu \in \mathfrak{h}^*\) . Show that \(\dim \operatorname{Hom}_{\mathfrak{g}}(\mathrm{Z}(\mu), \mathrm{Z}(\lambda)) \leq 1\) . (Let \(\varphi_1\) and \(\varphi_2\) be non-zero \(\mathfrak{g}\) -homomorphisms from \(\mathrm{Z}(\mu)\) to \(\mathrm{Z}(\lambda)\) . If \(\operatorname{Im}(\varphi_1) = \operatorname{Im}(\varphi_2)\) , \(\varphi_1\) and \(\varphi_2\) are linearly dependent by Prop. 1 (iii) of §6. Assume that \(\operatorname{Im}(\varphi_1) \neq \operatorname{Im}(\varphi_2)\) . If \(\mathrm{Z}(\mu)\) is simple, the sum \(\operatorname{Im}(\varphi_1) + \operatorname{Im}(\varphi_2)\) is direct; deduce that \(\mathfrak{P}(\xi + \lambda - \mu) \geq 2\mathfrak{P}(\xi)\) for all \(\xi \in \mathfrak{h}^*\) , contradicting \(b\) ). In the general case, use Exerc. 11.)
\end{enumerate}

When \(\dim\operatorname{Hom}_{\mathfrak{g}}(\mathrm{Z}(\mu),\mathrm{Z}(\lambda))=1\) , write \(\mathrm{Z}(\mu)\subset\mathrm{Z}(\lambda)\) by abuse of notation. d) Let \(\nu\in\mathfrak{h}^{*}\) . The set of \(\lambda\in\mathfrak{h}^{*}\) such that \(\mathrm{Z}(\lambda-\nu)\subset\mathrm{Z}(\lambda)\) is closed in \(\mathfrak{h}^{*}\) in the Zariski topology.

¶ 13) a) Let \(\mathfrak{a}\) be a nilpotent Lie algebra, \(x \in \mathfrak{a}, n \in \mathbf{N}, p \in \mathbf{N}\) . There exists \(l \in \mathbf{N}\) such that \(x^{l}y_{1} \ldots y_{n} \in \mathrm{U}(\mathfrak{a})x^{p}\) for all \(y_{1}, \ldots, y_{n} \in \mathfrak{a}\) .

\begin{enumerate}
\def\labelenumi{\alph{enumi})}
\setcounter{enumi}{1}
\tightlist
\item
  Let \(\lambda, \mu \in \mathfrak{h}^*\) , and \(\alpha \in \mathrm{B}\) be such that
\end{enumerate}

\[
\mathrm{Z} (s _ {\alpha} \mu - \rho) \subset \mathrm{Z} (\mu - \rho) \subset \mathrm{Z} (\lambda - \rho).
\]

Assume that \(\lambda \in \mathrm{P}\) . Let \(p = \lambda (H_{\alpha})\in \mathbf{Z}\) . Show that:

\(b_{1})\) If \(p \leq 0\) , then \(\mathrm{Z}(\lambda - \rho) \subset \mathrm{Z}(s_{\alpha}\lambda - \rho)\) .

\[
b _ {2}) \text {   If   } p > 0, \text {   then   } \mathrm{Z} (s _ {\alpha} \mu - \rho) \subset \mathrm{Z} (s _ {\alpha} \lambda - \rho) \subset \mathrm{Z} (\lambda - \rho).
\]

(Use \(a\) ), and §6, Cor. 1 of Prop. 6.)

\begin{enumerate}
\def\labelenumi{\alph{enumi})}
\setcounter{enumi}{2}
\tightlist
\item
  Let \(\lambda \in \mathfrak{h}^*\) , \(\alpha \in \mathrm{R}\) , and \(m = \lambda(H_{\alpha})\) . Assume that \(m \in \mathbf{N}\) . Show that
\end{enumerate}

\[
\mathrm{Z} (s _ {\alpha} \lambda - \rho) \subset \mathrm{Z} (\lambda - \rho).
\]

(Prove this first for \(\lambda \in \mathrm{P}\) by using \(b\) ), and then in the general case by using Exerc. 12 \(d\) .) \(^{17}\)

\begin{enumerate}
\def\labelenumi{\arabic{enumi})}
\setcounter{enumi}{13}
\item
  *Let \(\mathfrak{a}\) be a semi-simple Lie algebra, and \(Z(\mathfrak{a})\) the centre of \(U(\mathfrak{a})\) . Show that \(U(\mathfrak{a})\) is a free \(Z(\mathfrak{a})\) -module. (Remark that \(grU(\mathfrak{a})\) is isomorphic to \(S(\mathfrak{a})\) and \(S(\mathfrak{a}^*)\) , and use Remark 2 of no. 3.)*
\item
  Let \(x\) be a diagonalizable element of \(\mathfrak{g}\) (§3, Exerc. 10), and \(y\) a semi-simple element of \(\mathfrak{g}\) such that \(f(x) = f(y)\) for every invariant polynomial function \(f\) on \(\mathfrak{g}\) . Show that there exists \(s \in \mathrm{Aut}_e(\mathfrak{g})\) such that \(sy = x\) . (Remark that ad \(x\) and ad \(y\) have the same characteristic polynomial, cf.~Exerc. 8, hence the fact that y is diagonalizable. Next, reduce to the case in which x and y are contained in h by using Exerc. 10 of §3, and use Th. 1 (i) and Lemma 6 to prove that x and y are conjugate under W.)
\item
  Let \(\mathfrak{a}\) be a semi-simple Lie algebra, \(x\) an element of \(\mathfrak{a}\) , and \(x_{s}\) the semi-simple component of \(x\) . Show that, if \(f\) is an invariant polynomial function on \(\mathfrak{a}\) , then \(f(x) = f(x_{s})\) . (Reduce to the case in which \(f\) is of the form \(x \mapsto \operatorname{Tr} \rho(x)^n\) .)
\item
  Assume that \(k\) is algebraically closed, and put \(\mathrm{G} = \operatorname{Aut}_e(\mathfrak{g})\) . Let \(x, y \in \mathfrak{g}\) . Show the equivalence of:
\end{enumerate}

\begin{enumerate}
\def\labelenumi{(\roman{enumi})}
\item
  The semi-simple components of \(x\) and \(y\) are G-conjugate.
\item
  For every invariant polynomial function \(f\) on \(\mathfrak{g}\) , we have \(f(x) = f(y)\) . (Use Exerc. 15 and 16.)
\end{enumerate}

\begin{enumerate}
\def\labelenumi{\arabic{enumi})}
\setcounter{enumi}{17}
\tightlist
\item
  Let \(\mathfrak{a}\) be a semi-simple Lie algebra, \(l = \mathrm{rk}(\mathfrak{a})\) , I the algebra of invariant polynomial functions on \(\mathfrak{a}\) , and \(\mathrm{P}_1, \ldots, \mathrm{P}_l\) homogeneous elements of I generating I as an algebra. The \(\mathrm{P}_i\) define a polynomial map \(\mathrm{P} : \mathfrak{a} \to k^l\) . If \(x \in \mathfrak{a}\) , denote by \(\mathrm{D}_x\mathrm{P} : \mathfrak{a} \to k^l\) the tangent linear map of P at \(x\) (Chap. VII, App. I, no. 2).
\end{enumerate}

\begin{enumerate}
\def\labelenumi{\alph{enumi})}
\tightlist
\item
  Let \(\mathfrak{h}\) be a Cartan subalgebra of \(\mathfrak{a}\) , and let \(x \in \mathfrak{h}\) . Prove the equivalence of:
\end{enumerate}

\begin{enumerate}
\def\labelenumi{(\roman{enumi})}
\item
  \(D_{x}P|\mathfrak{h}\) is an isomorphism from \(\mathfrak{h}\) to \(k^{l}\) ;
\item
  \(x\) is regular.
\end{enumerate}

(Reduce to the split case. Choose a basis of \(\mathfrak{h}\) , and denote by \(d(x)\) the determinant of the matrix of \(\mathrm{D}_x\mathrm{P}|\mathfrak{h}\) relative to this basis. Show, by means of Prop. 5 of Chap. V, §5, no. 4, that there exists \(c \in k^*\) such that \(d(x)^2 = c \prod_{\alpha \in \mathrm{R}} \alpha(x)\) ,

where \(\alpha\) belongs to the set \(R\) of roots of \((\mathfrak{a},\mathfrak{h})\) .)

If these conditions are satisfied, \(\mathfrak{a} = \mathfrak{h} \oplus \operatorname{Im} \operatorname{ad}(x)\) , and \(\operatorname{Ker} D_x P = \operatorname{Im} \operatorname{ad}(x)\) (use Exerc. 6 to show that \(D_x P\) vanishes on \(\operatorname{Im} \operatorname{ad}(x)\) ).

\begin{enumerate}
\def\labelenumi{\alph{enumi})}
\setcounter{enumi}{1}
\tightlist
\item
  Show that the set of \(x \in a\) such that \(D_{x}P\) is of rank l is a dense open subset of a in the Zariski topology.
\end{enumerate}

\exercisesection{\S~9}

All the g-modules considered are assumed to be finite dimensional.

\begin{enumerate}
\def\labelenumi{\arabic{enumi})}
\item
  If \(m\) is an integer \(\geq 0\) , we have \(\dim \mathrm{E}(m\rho) = (m + 1)^{\mathrm{N}}\) , where \(\mathrm{N} = \operatorname{Card}(\mathrm{R}_{+})\) .
\item
  Show that there exists a unique polynomial function \(d\) on \(\mathfrak{h}^*\) such that \(d(\lambda) = \dim \mathrm{E}(\lambda)\) for all \(\lambda \in \mathrm{P}_{++}\) ; its degree is \(\operatorname{Card}(\mathrm{R}_+)\) . We have
\end{enumerate}

\[
d (w \lambda - \rho) = \varepsilon (w) d (\lambda - \rho) \quad \text { if } w \in W, \lambda \in \mathfrak {h} ^ {*}.
\]

In particular, the function \(\lambda \mapsto d(\lambda - \rho)^2\) is invariant under W. Deduce that there exists a unique element \(u\) of the centre of \(\mathrm{U}(\mathfrak{g})\) such that \(\chi_{\lambda}(u) = d(\lambda)^2\) for all \(\lambda \in \mathfrak{h}^*\) (apply Th. 2 of §8, no. 5). When \(\mathfrak{g} = \mathfrak{sl}(2,k)\) , we have \(u = C + 1\) , where \(C\) is the element defined in Exerc. 1 of §1.

\begin{enumerate}
\def\labelenumi{\arabic{enumi})}
\setcounter{enumi}{2}
\tightlist
\item
  Let \(k_{1}, \ldots, k_{l}\) be the characteristic degrees of the algebra of invariants of W (cf.~Chap. V, §5).
\end{enumerate}

\begin{enumerate}
\def\labelenumi{\alph{enumi})}
\item
  Show that, for all \(j \geq 1\) , the number of i such that \(k_{i} > j\) is equal to the number of \(\alpha \in R_{+}\) such that \(\langle \rho, H_{\alpha} \rangle = j\) . (Reduce to the case in which R is irreducible, and use Chap. VI, §4, Exerc. 6 c.) (Cf. §5, Exerc. 5 g).
\item
  Deduce the formula
\end{enumerate}

\[
\prod_ {\alpha \in \mathrm{R} _ {+}} \langle \rho , H _ {\alpha} \rangle = \prod_ {i = 1} ^ {l} (k _ {i} - 1)!.
\]

¶4) Assume that g is simple, and denote by \(\gamma\) the element of \(R_{+}\) such that \(H_{\gamma}\) is the highest root of \(R^{\vee}\) ; write \(H_{\gamma} = \sum_{\alpha \in B} n_{\alpha} H_{\alpha}\) . We have \(\langle \rho, H_{\gamma} \rangle = \sum n_{\alpha} = h - 1\) , where h is the Coxeter number of R (Chap. VI, §1, no. 11, Prop. 31).

\begin{enumerate}
\def\labelenumi{\alph{enumi})}
\tightlist
\item
  Let \(\alpha \in \mathrm{B}\) . Show that, for all \(\beta \in \mathrm{R}_{+}\) ,
\end{enumerate}

\[
\langle \varpi_ {\alpha} + \rho , H _ {\beta} \rangle \leq h + n _ {\alpha} - 1,
\]

and that equality holds if \(\beta = \gamma\) . Deduce that every prime factor of \(\dim \mathrm{E}(\varpi_{\alpha})\) is \(\leq h + n_{\alpha} - 1\) .

\begin{enumerate}
\def\labelenumi{\alph{enumi})}
\setcounter{enumi}{1}
\tightlist
\item
  Assume that \(\varpi_{\alpha}\) is not minuscule, i.e.~\(n_{\alpha} \geq 2\) . Let \(m \in [2, n_{\alpha}]\) and \(p = h + m - 1\) . Verify (cf.~Chap. VI, Plates) that there exists \(\beta \in \mathrm{R}_{+}\) such that \(\langle \varpi_{\alpha}, H_{\beta} \rangle = n_{\alpha}\) and \(\langle \rho, H_{\beta} \rangle = h - 1 - (n_{\alpha} - m)\) , hence \(\langle \varpi_{\alpha} + \rho, H_{\beta} \rangle = p\) . Deduce that, if \(p\) is prime, \(p\) divides \(\dim \mathrm{E}(\varpi_{\alpha})\) . (Remark that \(p\) does not divide any of the \(\langle \rho, H_{\beta} \rangle\) for \(\beta \in \mathrm{R}_{+}\) , cf.~Exerc. 3.)
\end{enumerate}

\(c\) ) When \(\mathfrak{g}\) is of type \(\mathrm{G}_2\) (resp. \(\mathrm{F}_4,\mathrm{E}_8\) ), we have \(h = 6\) (resp. 12, 30), and \(\dim \mathrm{E}(\varpi_{\alpha})\) is divisible by 7 (resp. 13, 31), When \(\mathfrak{g}\) is of type \(\mathrm{E}_6\) (resp. \(\mathrm{E}_7\) ), and \(\varpi_{\alpha}\) is not minuscule, \(\dim \mathrm{E}(\varpi_{\alpha})\) is divisible by 13 (resp. 19).

¶ 5) a) Let \(\alpha \in \mathrm{R}, x \in \mathfrak{g}^{\alpha}, y \in \mathfrak{g}^{-\alpha}\) , and let E be a \(\mathfrak{g}\) -module. Show that, for all \(\lambda \in \mathrm{P}\) ,

\[
\operatorname{Tr} ((x y) _ {\mathrm{E}} | \mathrm{E} ^ {\lambda}) = \operatorname{Tr} ((x y) _ {\mathrm{E}} | \mathrm{E} ^ {\lambda + \alpha}) + \lambda ([ x, y ]) \dim \mathrm{E} ^ {\lambda}.
\]

Deduce that

\[
\sum_ {\lambda \in P} \lambda ([ x, y ]) \dim E ^ {\lambda}. e ^ {\lambda} = (1 - e ^ {- \alpha}) \sum_ {\lambda \in P} \operatorname{Tr} ((x y) _ {E} | E ^ {\lambda}). e ^ {\lambda}.
\]

\begin{enumerate}
\def\labelenumi{\alph{enumi})}
\setcounter{enumi}{1}
\tightlist
\item
  Give \(\mathfrak{h}^*\) a non-degenerate W-invariant symmetric bilinear form \(\langle ,\rangle\) . Let \(\Delta\) be the endomorphism of the vector space \(k[\mathrm{P}]\) such that \(\Delta(e^{\mu}) = \langle \mu ,\mu \rangle e^{\mu}\) for all \(\mu \in \mathrm{P}\) ; if \(a,b\in k[\mathrm{P}]\) , put
\end{enumerate}

\[
\Delta^ {\prime} (a, b) = \Delta (a b) - a \Delta (b) - b \Delta (a).
\]

Prove that

\[
\Delta (\mathrm{J} (e ^ {\mu})) = \langle \mu , \mu \rangle \mathrm{J} (e ^ {\mu}) \quad \text { for } \mu \in \mathrm{P},
\]

\[
\Delta^ {\prime} (e ^ {\lambda}, e ^ {\mu}) = 2 \langle \lambda , \mu \rangle e ^ {\lambda + \mu} \quad \mathrm{for} \lambda , \mu \in \mathrm{P},
\]

\[
\Delta^ {\prime} (a b, c) = a \Delta^ {\prime} (b, c) + b \Delta^ {\prime} (a, c) \quad \text { for } a, b, c \in k [ P ].
\]

\(c)\) Let \(\lambda \in \mathrm{P}_{++}\) , \(c_{\lambda} = \operatorname{ch}(\operatorname{E}(\lambda))\) and \(d = \mathrm{J}(e^{\rho})\) . Prove that

\[
\Delta (c _ {\lambda} d) = \langle \lambda + \rho , \lambda + \rho \rangle c _ {\lambda} d.
\]

(Use \(a\) ), \(b\) ), §6, Cor. of Prop. 7, and Chap. VI, §3, no. 3, formula (3).)

\(d)\) Deduce another proof of the formula of H. Weyl from the above and §7, no. 2, Prop. 5 (iii).

\begin{enumerate}
\def\labelenumi{\alph{enumi})}
\setcounter{enumi}{4}
\tightlist
\item
  For all \(\lambda \in \mathfrak{h}^*\) , put \(\dim \mathrm{E}^{\lambda} = m(\lambda)\) . Deduce from \(a\) that
\end{enumerate}

\[
\begin{array}{c} \operatorname{Tr} ((x y) _ {\mathrm{E}} | \mathrm{E} ^ {\lambda}) = \sum_ {i = 0} ^ {+ \infty} (\lambda + i \alpha) ([ x, y ]) m (\lambda + i \alpha) \\ \sum_ {i = - \infty} ^ {+ \infty} (\lambda + i \alpha) ([ x, y ]) m (\lambda + i \alpha) = 0. \end{array}
\]

\(f)\) Let \(\langle\cdot,\cdot\rangle\) be a non-degenerate invariant symmetric bilinear form on \(\mathfrak{g}\) whose restriction to \(\mathfrak{h}\) is the inverse of the form chosen above. Let \(\Gamma\) be the corresponding Casimir element. Assume that E is simple; put \(\Gamma_{\mathrm{E}}=\gamma.1\) , where \(\gamma\in k\) . By using \(e\) and §2, no. 3, Prop. 6, show that

\[
\gamma m (\lambda) = \langle \lambda , \lambda \rangle m (\lambda) + \sum_ {\alpha \in \mathrm{R}} \sum_ {i = 0} ^ {+ \infty} \langle \lambda + i \alpha , \alpha \rangle m (\lambda + i \alpha)
\]

for all \(\lambda \in \mathfrak{h}^*\) , and then that

\[
\begin{array}{l} \gamma m (\lambda) = \langle \lambda , \lambda \rangle m (\lambda) + \sum_ {\alpha \in \mathrm{R} _ {+}} m (\lambda) \langle \lambda , \alpha \rangle + 2 \sum_ {\alpha \in \mathrm{R} _ {+}} \sum_ {i = 1} ^ {+ \infty} m (\lambda + i \alpha) \langle \lambda + i \alpha , \alpha \rangle \\ = \langle \lambda , \lambda + 2 \rho \rangle m (\lambda) + 2 \sum_ {\alpha \in \mathrm{R} _ {+}} \sum_ {i = 1} ^ {+ \infty} m (\lambda + i \alpha) \langle \lambda + i \alpha , \alpha \rangle . \end{array}
\]

\(g)\) Continue to assume that E is simple; let \(\omega\) be its highest weight. Deduce from \(f)\) that, for all \(\lambda \in \mathfrak{h}^*\) ,

\[
(\langle \omega + \rho , \omega + \rho \rangle - \langle \lambda + \rho , \lambda + \rho \rangle) m (\lambda) = 2 \sum_ {\alpha \in \mathrm{R} _ {+}} \sum_ {i = 1} ^ {+ \infty} m (\lambda + i \alpha) \langle \lambda + i \alpha , \alpha \rangle .
\]

(Recall that, by Prop. 5 of §7, \(\langle\omega+\rho,\omega+\rho\rangle>\langle\lambda+\rho,\lambda+\rho\rangle\) if \(\lambda\) is a weight of E distinct from \(\omega\) . The preceding formula thus gives a procedure for calculating \(m(\lambda)\) step by step.) \(^{18}\)

\begin{enumerate}
\def\labelenumi{\arabic{enumi})}
\setcounter{enumi}{5}
\tightlist
\item
  Let \(x \mapsto x^*\) be the involution of \(k[\mathrm{P}]\) that takes \(e^p\) to \(e^{-p}\) for all \(p \in \mathrm{P}\) . \(a)\) Put \(\mathrm{D} = d^* d = \prod_{\alpha \in \mathrm{R}} (1 - e^\alpha)\) . Show that \(d^* = (-1)^{\mathrm{N}} d\) , where \(\mathrm{N} = \mathrm{Card}(\mathrm{R}_+)\) , and hence that \(\mathrm{D} = (-1)^{\mathrm{N}} d^2\) .
\end{enumerate}

\begin{enumerate}
\def\labelenumi{\alph{enumi})}
\setcounter{enumi}{1}
\tightlist
\item
  Define two linear forms \(\varepsilon\) and I on \(k[\mathrm{P}]\) by the formulas:
\end{enumerate}

\[
\varepsilon (1) = 1, \quad \varepsilon (e ^ {p}) = 0 \quad \text { if } p \in \mathrm{P} - \{0 \}
\]

and

\[
\mathrm{I} (f) = \frac {1}{m} \varepsilon (\mathrm{D}. f), \quad \mathrm{where} m = \mathrm{Card} (\mathrm{W}).
\]

Show, by using the formula \(d = \mathrm{J}(e^{\rho})\) , that \(\mathrm{I}(1) = 1\) .

\begin{enumerate}
\def\labelenumi{\alph{enumi})}
\setcounter{enumi}{2}
\tightlist
\item
  Let \(\lambda \in \mathrm{P}_{++}\) and \(c_{\lambda} = \operatorname{ch} \mathrm{E}(\lambda) = \mathrm{J}(e^{\lambda + \rho}) / d\) . Show that \(\mathrm{I}(c_{\lambda}) = 0\) if \(\lambda \neq 0\) . (Same method as for \(b\) .)
\end{enumerate}

\(d\) ) Show that I takes integer values on the subalgebra \(\mathbf{Z}[\mathrm{P}]^{\mathrm{W}} = \operatorname {ch}\operatorname {R}(\mathfrak{g})\) of \(k[\mathrm{P}]\) . If E is a \(\mathfrak{g}\) -module, the dimension of the space of invariants of \(\mathfrak{g}\) in E is equal to I(ch E). (Reduce to the case in which E is simple and use \(b\) ) and \(c\) .) \(e\) ) We have \(\dim \mathrm{E} = \sum_{\lambda \in \mathrm{P}_{++}}\mathrm{I}(c_\lambda^*\mathrm{ch}\mathrm{E})d(\lambda)\) , where \(d(\lambda) = \dim \mathrm{E}(\lambda)\) . In particular:

\[
\mathrm{I} (c _ {\lambda} ^ {*} c _ {\mu}) = \delta_ {\lambda \mu} \quad \text { if } \lambda , \mu \in \mathrm{P} _ {+ +}.
\]

\(f)\) If \(\lambda, \mu, \nu \in \mathrm{P}_{++}\) , the integer \(m(\lambda, \mu, \nu)\) of Prop. 2 is equal to \(\mathrm{I}(c_{\lambda} c_{\mu} c_{\nu}^{*})\) . Deduce the identity

\[
d (\lambda) d (\mu) = \sum_ {\nu \in \mathrm{P} _ {+ +}} m (\lambda , \mu , \nu) d (\nu) \quad \lambda , \mu , \nu \in \mathrm{P} _ {+ +}.
\]

(Apply e) to the g-module \(\mathrm{E} = \mathrm{E}(\lambda) \otimes \mathrm{E}(\mu).\)

¶ 7) We retain the notations of the proof of Th. 2.

\begin{enumerate}
\def\labelenumi{\alph{enumi})}
\tightlist
\item
  Show that
\end{enumerate}

\[
f _ {\rho} (\mathrm{J} (e ^ {\mu})) = \prod_ {\alpha \in \mathrm{R} _ {+}} (e ^ {(\mu | \alpha) \mathrm{T} / 2} - e ^ {- (\mu | \alpha) \mathrm{T} / 2}).
\]

\begin{enumerate}
\def\labelenumi{\alph{enumi})}
\setcounter{enumi}{1}
\tightlist
\item
  Take \((\cdot |\cdot)\) to be the canonical bilinear form \(\varPhi_{\mathrm{R}}\) (Chap. VI, §1, no. 12). Show that
\end{enumerate}

\[
  f _ {\rho} (\mathrm{J} (e ^ {\mu})) \equiv d _ {\mu} \mathrm{T} ^ {\mathrm{N}} \left(1 + \frac {\mathrm{T} ^ {2}}{48} (\mu | \mu)\right) \pmod {\mathrm{T} ^ {\mathrm{N} + 3} \mathbf {R} [ [ \mathrm{T} ] ]},
\]

where \(d_{\mu} = \prod_{\alpha \in \mathrm{R}_{+}}(\mu |\alpha)\) .

\begin{enumerate}
\def\labelenumi{\alph{enumi})}
\setcounter{enumi}{2}
\tightlist
\item
  From b) and the equality \(\mathrm{J}(e^{\lambda+\rho})=\mathrm{ch}(\mathrm{E}).\mathrm{J}(e^{\rho})\) , deduce the formula
\end{enumerate}

\[
  \sum_ {\mu \in \mathrm{P}} (\mu | \rho) ^ {2} \dim \mathrm{E} ^ {\mu} = \frac {\dim \mathrm{E}}{24} (\lambda | \lambda + 2 \rho).
\]

\(d\) ) Assume that \(\mathfrak{g}\) is simple. Show that \((\rho |\rho) = \dim \mathfrak{g} / 24\) . (Apply \(c\) ) with \(\lambda\) the highest root of R, and use Exerc. 3 of §6.)

¶8) Let \(\psi\) be a polynomial function on \(\mathfrak{h}^*\) of degree \(r\) . Show that there exists a unique polynomial function \(\varPsi\) on \(\mathfrak{h}^*\) that is invariant under W, of degree \(\leq r\) , and such that

\[
\sum_ {\mu \in \mathrm{P}} \psi (\mu) \dim \mathrm{E} ^ {\mu} = \varPsi (\lambda + \rho) \dim \mathrm{E}
\]

for any simple g-module E of highest weight \(\lambda\) .

(Treat first the case in which \(\psi(\mu)=(\mu|\nu)^{r}\) , where \(\nu\in P\) is not orthogonal to any root; for this use the homomorphism \(f_{\nu}\) in the proof of Th. 2 and Chap. V, §5, no. 4, Prop. 5 (i).)

\begin{enumerate}
\def\labelenumi{\arabic{enumi})}
\setcounter{enumi}{8}
\tightlist
\item
  We use the notations of Chap. VI, Plate I, in the case of an algebra \(\mathfrak{g}\) of type \(\mathrm{A}_2\) . Let \(n, p\) be integers \(\geq 0\) .
\end{enumerate}

\[
a) \mathfrak {P} (n \alpha_ {1} + p \alpha_ {2}) = 1 + \inf (n, p).
\]

\begin{enumerate}
\def\labelenumi{\alph{enumi})}
\setcounter{enumi}{1}
\tightlist
\item
  Let \(\lambda = n\varpi_1 + p\varpi_2\) . Then \(\dim \mathrm{E}(\lambda) = \frac{1}{2}(n + 1)(p + 1)(n + p + 2)\) . The multiplicity of the weight 0 of \(\mathrm{E}(\lambda)\) is 0 if \(\lambda\) is not radical, i.e.~if \(n \not\equiv p \pmod{3}\) ; if \(\lambda\) is radical, it is \(1 + \inf(n,p)\) .
\end{enumerate}

\begin{enumerate}
\def\labelenumi{\arabic{enumi})}
\setcounter{enumi}{9}
\tightlist
\item
  We use the notations of Chap. VI, Plate II, in the case of an algebra of type \(B_{2}\) . Let n, p be integers \(\geq 0\) .
\end{enumerate}

\begin{enumerate}
\def\labelenumi{\alph{enumi})}
\tightlist
\item
  We have
\end{enumerate}

\[
\begin{array}{c} \mathfrak {P} (n \alpha_ {1} + p (\alpha_ {1} + 2 \alpha_ {2})) = 1 + \frac {1}{2} p (p + 3) \\ \mathfrak {P} (n \alpha_ {2} + p (\alpha_ {1} + \alpha_ {2})) = [ p ^ {2} / 4 ] + p + 1 \\ \mathfrak {P} (n (\alpha_ {1} + \alpha_ {2}) + p (\alpha_ {1} + 2 \alpha_ {2})) = [ n ^ {2} / 4 ] + n + 1 + n p + \frac {1}{2} p (p + 3). \end{array}
\]

\begin{enumerate}
\def\labelenumi{\alph{enumi})}
\setcounter{enumi}{1}
\tightlist
\item
  Let \(\lambda = n(\alpha_{1} + \alpha_{2}) + p(\alpha_{1} + 2\alpha_{2})\) . The multiplicity of the weight 0 in \(\operatorname{E}(\lambda)\) is
\end{enumerate}

\[
[ n / 2 ] + 1 + n p + p.
\]

\begin{enumerate}
\def\labelenumi{\arabic{enumi})}
\setcounter{enumi}{10}
\item
  Assume that \(\mathfrak{g}\) is not a product of algebras of rank 1. Let \(n\in \mathbf{N}\) . Show that there exists a simple \(\mathfrak{g}\) -module one of whose weights is of multiplicity \(\geq n\) . (Suppose not. Let \(E_{\lambda}\) be the simple g-module of highest weight \(\lambda \in P_{++}\) . Compare \(\dim E_{\lambda}\) and the number of distinct weights of \(E_{\lambda}\) when \((\lambda|\lambda) \to \infty\) (with the notations of Th. 2).)
\item
  Let \(U^{0}\) be the commutant of h in \(\mathrm{U}(\mathfrak{g})\) . If g is of rank 1, \(U^{0}\) is commutative. If g is of rank \(\geq 2\) , \(U^{0}\) admits simple representations of arbitrarily large finite dimension. (Use Exerc. 11 and Exerc. 23 a) of §7.)
\item
  Let R be a root system in a vector space V. Two elements \(v_{1}, v_{2}\) of V are said to be disjoint if R is the direct sum of two root systems \(R_{1}\) and \(R_{2}\) (Chap. VI, §1, no. 2) such that \(v_{i}\) belongs to the vector subspace of V generated by \(R_{i}, i = 1, 2\) . Show that two elements of \(V_{R}\) that belong to the same chamber of R are disjoint if and only if they are orthogonal.
\end{enumerate}

¶ 14) a) Let \(\mu, \nu \in \mathrm{P}_{++}\) and let \(\gamma\) be a weight of \(\mathrm{E}(\mu)\) . Let \(\rho_{\mu}\) be the representation of \(\mathfrak{g}\) on \(\mathrm{E}(\mu)\) . Let \(X_{\alpha} \in \mathfrak{g}^{\alpha} - \{0\}, Y_{\alpha} \in \mathfrak{g}^{-\alpha} - \{0\}\) . If \(\alpha \in \mathrm{B}\) , put \(\nu_{\alpha} = \nu(H_{\alpha})\) , and

\[
\mathrm{E} ^ {+} (\mu , \gamma , \nu) = \mathrm{E} (\mu) ^ {\gamma} \cap \bigcap_ {\alpha \in \mathrm{B}} \mathrm{Ker} \rho_ {\mu} (X _ {\alpha}) ^ {\nu_ {\alpha} + 1},
\]

\[
\mathrm{E} ^ {-} (\mu , \gamma , \nu) = \mathrm{E} (\mu) ^ {\gamma} \cap \bigcap_ {\alpha \in \mathrm{B}} \operatorname{Ker} \rho_ {\mu} (Y _ {\alpha}) ^ {\nu_ {\alpha} + 1},
\]